\documentclass[11pt]{article}
\usepackage[a4paper,margin=22mm]{geometry}
\usepackage[no-math]{fontspec}
\setmainfont{lmroman10-regular.otf}[BoldFont=lmroman10-bold.otf,ItalicFont=lmroman10-italic.otf,BoldItalicFont=lmroman10-bolditalic.otf]
\usepackage{amsmath,amssymb,amsthm,mathtools,graphicx,xcolor,hyperref,xurl,xspace,ifthen}
\setlength{\emergencystretch}{4em}
\SetMathAlphabet{\mathbf}{normal}{OT1}{cmr}{bx}{n}
\SetMathAlphabet{\mathbf}{bold}{OT1}{cmr}{bx}{n}
\newtheorem{theo}{Theorem}[section]
\newtheorem{lemm}[theo]{Lemma}
\newtheorem{thm}[theo]{Theorem}
\newtheorem{theorem}[theo]{Theorem}
\newtheorem*{remas*}{Remarks}
\newtheorem{remas}[theo]{Remarks}
\newtheorem*{exam*}{Example}
\newtheorem{ques}[theo]{Question}
\providecommand{\og}{«\nobreakspace}
\providecommand{\fg}{\nobreakspace»}
\newtheorem{lem}[theo]{Lemma}
\newtheorem{lemma}[theo]{Lemma}
\newtheorem{prop}[theo]{Proposition}
\newtheorem{coro}[theo]{Corollary}
\newtheorem{corol}[theo]{Corollary}
\newtheorem{conj}[theo]{Conjecture}
\newtheorem{Proposition}[theo]{Proposition}
\newtheorem{theoreme}[theo]{Theorem}
\theoremstyle{definition}
\newtheorem{defi}[theo]{Definition}
\newtheorem{exam}[theo]{Example}
\newtheorem{exem}[theo]{Example}
\newtheorem{remark}[theo]{Remark}
\newtheorem{example}[theo]{Example}
\newtheorem{rema}[theo]{Remark}
\newtheorem{nota}[theo]{Notation}
\newtheorem*{rema*}{Remark}
\newtheorem*{defi*}{Definition}
\newtheorem*{theo*}{Theorem}
\newtheorem*{lemm*}{Lemma}
\newtheorem*{prop*}{Proposition}
\newcommand{\enoncename}{}
\newtheorem{corpusenonce}[theo]{\enoncename}
\newenvironment{enonce}[2][]{\renewcommand{\enoncename}{#2}\begin{corpusenonce}}{\end{corpusenonce}}
\newenvironment{enonce*}[2][]{\par\medskip\noindent\textbf{#2.}\itshape}{\par\medskip}
\newenvironment{proofc}[1][\proofname]{\begin{proof}[#1]}{\end{proof}}
\newcommand{\Mc}{\mathcal{M}}
\newcommand{\pointir}{.\ ---\ }
\providecommand{\mathof}[1]{\mathrm{#1}}

\numberwithin{equation}{section}
\numberwithin{figure}{section}
\numberwithin{table}{section}
\def\endappendix{}
%~Mouliné par MaN_auto v.0.21.2 2019-02-11 15:07:16


\usepackage[all]{xy}







\DeclareMathOperator{\im}{Im}
\DeclareMathOperator{\Id}{Id}
%\DeclareMathOperator{\vol}{vol}
\DeclareMathOperator{\Trace}{Trace}
\DeclareMathOperator{\trace}{Trace}

\newcommand{\R}{\mathbb{R}}
\newcommand{\N}{\mathbb{N}}

\newcommand*\eps{\epsilon}
\newcommand*\calS{\mathcal{S}}
\newcommand*\F{\mathcal{F}}
\newcommand*\Fi{\mathcal{F}^{-1}}

\newcommand{\dd}{\mathrm{d}}
\newcommand{\dt}{\dd t}
\newcommand{\ds}{\dd s}
\newcommand{\du}{\dd u}
\newcommand{\dv}{\dd v}
\newcommand{\dx}{\dd x}
\newcommand{\dK}{\dd K}
\newcommand{\dphi}{\dd \phi}
\newcommand{\drho}{\dd \rho}
\newcommand{\dalpha}{\dd \alpha}
\newcommand{\domega}{\dd \omega}
\newcommand{\dvol}{\dd \mathrm{vol}}


\graphicspath{{./figures/}}

\newcommand*{\mk}{\mkern -1mu}
\newcommand*{\Mk}{\mkern -2mu}
\newcommand*{\mK}{\mkern 1mu}
\newcommand*{\MK}{\mkern 2mu}

\hypersetup{urlcolor=purple, linkcolor=blue, citecolor=red}

\newcommand*{\romanenumi}{\renewcommand*{\theenumi}{\roman{enumi}}}
\newcommand*{\Romanenumi}{\renewcommand*{\theenumi}{\Roman{enumi}}}
\newcommand*{\alphenumi}{\renewcommand*{\theenumi}{\alph{enumi}}}
\newcommand*{\Alphenumi}{\renewcommand*{\theenumi}{\Alph{enumi}}}


\begin{document}
\section*{\raggedright 2749: Vanishing of finite-exponent ell-q-p cohomology of Carnot groups above the maximal consecutive Lie-cohomology weight gap}
\noindent\textbf{Author:} Pierre Pansu; Michel Rumin\par\noindent\textbf{Source:} On the \(\ell\)\textasciicircum{}\{q,p\} cohomology of Carnot groups. Annales Henri Lebesgue 1 (2018), publisher version, DOI 10.5802/ahl.8. \url{https://ahl.centre-mersenne.org/item/AHL_2018__1__267_0/}
\par\noindent\textbf{License:} CC BY 4.0, \url{https://creativecommons.org/licenses/by/4.0/}. The original publisher PDF has the explicit license notice. See the saved source/license records.
\par\noindent\textbf{Editorial material note (not author text):} The selected problem is the vanishing part (1) of Theorem 1.3 for a Carnot group, finite real exponents 1<p,q<infinity and the exact maximal consecutive Lie-cohomology weight-gap threshold. The complete author body retains all local Poincare, discrete/Sobolev equivalence, graded differential homotopy, all-moments-zero density/extension, integrability and weight-gap arguments. Explicitly published contracted-complex and pseudodifferential inversion results remain named prerequisites; nonvanishing and Engel-example assertions are unselected context. Original author discussion headings and prose are preserved. All proof text and formulas below are editable author TeX. Mathematical wording, language and proof order are retained. Only the article class, title matter and bibliography presentation are adapted for independent compilation; no proof has been generated or repaired.
\par\noindent\textbf{Editorial difficulty:} H2: The proof constructs a discrete-to-Sobolev cohomology equivalence, extends a graded pseudodifferential homotopy via an all-moments-zero density construction, and converts Lie-cohomology weights into global integrability gains. This nonelliptic cohomological analytic mechanism exceeds ordinary qualifying geometry and analysis.
\par\noindent\textbf{Editorial retrieval and attribution:} Retrieved 2026-10-01. DOI 10.5802/ahl.8. Source attribution notice: ../licenses/2749-source-license.txt; complete legal text: ../licenses/CC-BY-4.0.txt. Original source excerpts and the existing source-material check are retained in ../sources/2749/.
\medskip\hrule\medskip
\noindent\textbf{Author statement, complete proof and local supporting text follow in source order.}
\section{Introduction}

\subsection{\texorpdfstring{$\ell^{q,p}$}{ell(q,p)} cohomology}

Let $T$ be a countable simplicial complex. Given $1\leq p\leq q\leq +\infty$, the $\ell^{q,p}$ cohomology of $T$ is the quotient of the space of $\ell^p$ simplicial cocycles by the image of $\ell^q$ simplicial cochains by the coboundary $\dd$,
\[
\ell^{q,p}H^k(T)=(\ell^p C^k(T)\cap \ker \dd)/\dd(\ell^q C^{k-1}(T)) \cap \ell^p C^k(T).
\]
It is a quasiisometry invariant of bounded geometry simplicial complexes whose usual cohomology vanishes in a uniform manner, see~\cite{Bourdon-Pajot,Elek,Genton,Gromov,Pansu2}. Riemannian manifolds $M$ with bounded geometry admit quasiisometrically homotopy equivalent simplicial complexes (a construction is provided below, in Section~\ref{sec:leray}). Uniform vanishing of cohomology passes through. Therefore one can take the $\ell^{q,p}$ cohomology of any such complex as a definition of the $\ell^{q,p}$ cohomology of $M$.

$\ell^{q,p}$ cohomology has an analytic flavor. Its vanishing is equivalent to a \emph{discrete Poincar\'e inequality}: every $\ell^p$ $k$-cocycle $\kappa$ admits an $\ell^q$ primitive, i.e. a $(k-1)$-cochain $\lambda$ such that $\dd\lambda=\kappa$ and
\begin{equation}\label{eq:1}
\|\lambda\|_q \leq C \|\kappa\|_p.
\end{equation}

However, one may also think of $\ell^{q,p}$ cohomology as a (large scale) topological invariant. It has been useful in several contexts, mainly for the class of hyperbolic groups where the relevant value of $q$ is $q=p$, see~\cite{Bourdon-Kleiner,Bourdon-Pajot, Cornulier-Tessera,Drutu-Mackay} for instance. For this class of groups, $\ell^{p,p}$ cohomology is often Hausdorff and can sometimes be viewed as a function space of objects (functions, forms,\dots) living on the ideal boundary, opening the way to analysis on the ideal boundary. The intervals where $\ell^{p,p}$ cohomology vanishes (resp. is Hausdorff) provide numerical quasiisometry invariants of hyperbolic groups and negatively curved manifolds. In Riemannian geometry, negative curvature pinching seems to have a direct relation with $\ell^{p,p}$ cohomology, see~\cite{Pansu4}.

It is interesting to study a class of spaces where values of $q\not=p$ play a significant role. The goal of the present paper is to compute $\ell^{q,p}$ cohomology, to some extent, for certain Carnot groups. Even the case of abelian groups is not straightforward. A more remote goal would be to cover the whole family of groups of polynomial growth, the quasiisometry classification of which is still open. In this class, $\ell^{q,p}$ cohomology is expected to be either $0$ or nonHausdorff. Indeed, a result of Gromov, \cite[8.C]{Gromov}, states that for $p>1$, the reduced $\ell^{p,p}$ cohomology of a group with infinite center must vanish. Nevertheless, the intervals of $p$ and $q$ where $\ell^{q,p}$ cohomology vanishes could be exploited. Unfortunately, our present results have no direct consequences for the quasiisometry classification problem, since the quasiisometry classification of Carnot groups is known.

Before stating a general result, we shall describe a few examples, and indicate ideas along the way. A more thorough explanation of the methods used appears in Section~\ref{sec:method}.

\subsection{The Euclidean Sobolev inequality}

In Euclidean space $\R^n$, the Sobolev inequality states that there exists a constant $C=C(n)$ such that for every smooth compactly supported function $u$,
\begin{equation}\label{eq:2}
\|u\|_q \leq C \|\du\|_p,
\end{equation}
provided $p\geq 1$ and
\begin{equation*}
\frac{1}{p}-\frac{1}{q}=\frac{1}{n}.
\end{equation*}

Here is a sophisticated proof, valid if $p>1$. Let $\dd^*$ denote the adjoint of $\dd$, a differential operator from $1$-forms to functions. Let $\Delta=\dd^*\dd$ be the Laplacian. Then $\Delta$ admits a pseudodifferential inverse $\Delta^{-1}$ which commutes with $\dd$. Let $K=\Delta^{-1}\dd^*$. Then $\dK=1$, hence $v=K\du$ satisfies $\du=\dv$. Calderon--Zygmund theory shows that if $p>1$ and $\frac{1}{p}-\frac{1}{q}=\frac{1}{n}$, $K$ is bounded from $L^p$ to $L^q$. Hence $v=u$ and $\|u\|_q=\|K\du\|_q\leq C\,\|\du\|_p$.

This argument proves a slightly different statement, the \emph{continuous Poincar\'e inequality}: under the same assumption on $p>1$ and $q$, there exists a constant $C=C(n)$ such that for every smooth function $u$ such that $\du\in L^p$, there exists $c_u\in\R$ such that
\begin{equation}\label{eq:3}
\|u-c_u\|_q \leq C \|\du\|_p.
\end{equation}
This is strongly reminiscent of the discrete Poincar\'e inequality~\eqref{eq:1} in degree $k=1$.

Conversely, if $\frac{1}{p}-\frac{1}{q}\not=\frac{1}{n}$, sequences of functions yielding counterexamples to~\eqref{eq:2} can be constructed from functions which are homogeneous under the $1$-parameter group of homotheties $s\mapsto h_s=e^{s}\Id$ of $\R^n$. If $\frac{1}{p}-\frac{1}{q}>\frac{1}{n}$, examples exist with compact support in a fixed ball: the inequality fails for local reasons. If $\frac{1}{p}-\frac{1}{q}<\frac{1}{n}$, examples exist which satisfy a uniform bound on any finite number of derivatives: the inequality fails for large scale reasons. This dichotomy indicates that~\eqref{eq:2} is not equivalent to~\eqref{eq:1}. The continuous analogue of~\eqref{eq:1} amounts to requiring that~\eqref{eq:2} holds with Sobolev norms involving many derivatives instead of mere $L^p$ norms, see Theorem~\ref{prop:leray}.

Once the discretization dictionary is established, validity of~\eqref{eq:2} for $\frac{1}{p}-\frac{1}{q}=\frac{1}{n}$ implies validity of~\eqref{eq:1}. Since $\ell^{q'}\subset \ell^{q}$ for $q'\leq q$, the validity of~\eqref{eq:1} for $\frac{1}{p}-\frac{1}{q}\geq\frac{1}{n}$ follows automatically.

\subsection{Higher degree forms on abelian groups}

As we have just seen, inequality~\eqref{eq:2} is related to cohomology in degree~$1$. In higher degrees, let us define the continuous Poincar\'e inequality as follows: there exists a constant $C$ such that every closed differential $k$-form $\omega$ admits a primitive, i.e. a differential $k-1$-form $\phi$ such that $\dphi=\omega$ and
\begin{equation}\label{eq:4}
\|\phi\|_q \leq C \|\omega\|_p,
\end{equation}
where $\dd$ now denotes the exterior differential. For which exponents does such an inequality hold? Homogeneity under homotheties imposes a restriction. If $\omega$ is a differential $k$-form on $\R^n$, then
\begin{equation}\label{eq:5}
\|h_s^*\omega\|_p=e^{s(\frac{n}{p}-k)}\|\omega\|_p.
\end{equation}
Therefore in order for inequality~\eqref{eq:4} to hold for every closed $k$-form $\omega$ on $\R^n$, one needs that $\frac{n}{p}-k=\frac{n}{q}-(k-1)$, i.e.
\begin{equation*}
\frac{1}{p}-\frac{1}{q}=\frac{1}{n},
\end{equation*}
like in degree~$1$.

The argument based on properties of the scalar Laplacian extends to other degrees. To prove~\eqref{eq:4}, a homotopy $K$ is constructed. This is an operator on differential forms that decreases the degree by $1$ unit and satisfies $1=\dK+K\dd$. If $\dd\omega=0$, then $\omega=\dd(K\omega)$. Our favorite choice is $K=\dd^*\Delta^{-1}$, where $\dd^*$ is the $L^2$ adjoint of $\dd$, $\Delta= \dd\dd^*+\dd^*\dd$ and $\Delta^{-1}$ is a pseudodifferential inverse of $\Delta$, which commutes with $\dd$. Calderon--Zygmund theory shows again that if $p>1$ and $\frac{1}{p}-\frac{1}{q}=\frac{1}{n}$ is bounded from $L^p$ to $L^q$.

We shall see that this argument can be adapted to Carnot groups, with substantial changes.

\subsection{The Carnot Sobolev inequality}

A \emph{homogeneous group} is a Lie group $G$ quipped with a $1$-parameter group $s\mapsto h_s$ of expanding automorphisms. Expanding means that the eigenvalues of the derivation $D$ generating the $1$-parameter group are real and positive. When the eigenspace relative to eigenvalue $1$ generates the Lie algebra $\mathfrak{g}$, one speaks of a \emph{Carnot group}.

The derivation defines gradations, called \emph{weight} gradations, on $\mathfrak{g}$ and $\Lambda^\cdot\mathfrak{g}^*$. The trace of $D$,
\begin{equation*}
Q=\Trace(D)
\end{equation*}
is called the \emph{homogeneous dimension} of $(G,\{h_s\})$.

\begin{exem}
For abelian groups $G$, the derivation $D$ is the identity, all $k$-forms have weight $k$, $Q=\dim(G)$.

For Heisenberg groups, $D$ has two eigenvalues, $1$ of even multiplicity and $2$ of multiplicity $1$. For $1\leq k\leq \dim(G)-1$, $\Lambda^k\mathfrak{g}^*$ splits into two eigenspaces with weights $k$ and $k+1$, $Q=\dim(G)+1$.
\end{exem}

When $G$ is not abelian, $1$-forms come in several different weights $1,2,\dots,r$ where $r$ is the step of $G$. Assume that an inequality of the form~\eqref{eq:4},
\begin{equation*}
\|u\|_q \leq C \|\du\|_p
\end{equation*}
holds for all compactly supported functions $u$ on $G$. Splitting $\du=\dd_1 u+\dots +\dd_r u$ into weights and applying the obvious generalization of Formula~\eqref{eq:5}, we get, for every $s\in\R$,
\begin{equation*}
e^{s\frac{Q}{q}} \|u\|_q \leq C' \left(e^{s(\frac{Q}{p}-1)}\|\dd_1 u\|_p+\dots +e^{s(\frac{Q}{p}-r)}\|\dd_r u\|_p\right).
\end{equation*}
Letting $s$ tend to $+\infty$, we find two necessary conditions:
\begin{equation}\label{eq:6}
\frac{1}{p}-\frac{1}{q}=\frac{1}{Q} \quad \text{and}\quad \|u\|_q \leq C' \|\dd_1 u\|_p.
\end{equation}
It turns out that such an inequality indeed holds, see~\cite{FollandStein} for $p>1$ and~\cite{VCSC} for $p=1$.

\subsection{Higher degree forms on Carnot groups}

For higher degrees, one can similarly reduce the range of weights occurring in each degree. The \emph{contracted complex} is a resolution of the map $\dd_1$ from $0$ to $1$-forms along $\mathfrak{g}_1$. It is homotopy equivalent to de Rham's complex. The weights of forms occurring in the contracted complex are the same which occur in Lie algebra cohomology. This is an important step: if, in two consecutive degrees, Lie algebra cohomology exists in only one weight, then the differential of the contracted complex is homogeneous, and a variant of the method based on inverting the Laplacian leads to sharp intervals for vanishing and nonvanishing of $\ell^{q,p}$-cohomology. This happens in all degrees for Heisenberg groups, for instance. This approach has been pursued in~\cite{BFP}.

In this paper, we shall obtain results in the general case, using a different method. The main trick, stated in~\cite{Rumin_CRAS_1999}, is a homogeneous (pseudodifferential) version of the exterior differential. Inverting the corresponding Laplacian leads to a homotopy $K$ which is bounded on certain function spaces of Sobolev type where each weight component of a form is differentiated as many times as its weight. A Poincar\'e inequality for functions, similar to~\eqref{eq:6}, allows to get back to usual Lebesgue spaces, but with losses caused by lack of homogeneity. More details on the method will be given after the statement of our main theorem.

\subsection{Main result}

Our main result relates the vanishing and non-vanishing of $\ell^{q,p}$ cohomology of a Carnot group $G$ to the weight gaps in its Lie algebra cohomology $H^*(\mathfrak{g})$ in two consecutive degrees. Recall that this cohomology can be seen as the cohomology of translation invariant forms on $G$. It is graded by degree and weight,
\[
H^\cdot(\mathfrak{g})=\bigoplus_{k,w}H^{k,w}(\mathfrak{g}).
\]

\begin{defi}
For $k=0,\dots,\dim(\mathfrak{g})$, let $w_{\min}(k)$ (resp. $w_{\max}(k)$) be the smallest (resp. the largest) weight $w$ such that $H^{k,w}(\mathfrak{g})\not=0$.
\end{defi}

For instance, on $G=\R^n$, all invariant forms are closed, therefore $H^*(\mathfrak{g}) = \Lambda^* \mathfrak{g}$, and $w_{\min}(k) =w_{\max}(k) = k$ for $k=0, \dots, n$.

\begin{theo}\label{1}
Let $G$ be a Carnot group of dimension $n$ and of homogeneous dimension $Q$. Let $k=1,\dots,n$. Denote by
\[
\delta N_{\max}(k)=w_{\max}(k)-w_{\min}(k-1),\quad \delta N_{\min}(k)=\max\{1,w_{\min}(k)-w_{\max}(k-1)\}.
\]
Let $p$ and $q$ be real numbers.
\begin{enumerate}%[label=(\roman*)]
\item \label{1.3.1} If
\[
1<p,q<\infty \quad \text{and}\quad
\frac{1}{p}-\frac{1}{q}\geq\frac{\delta N_{\max}(k)}{Q}.
\]
Then the $\ell^{q,p}$ cohomology in degree $k$ of $G$ vanishes.
\item \label{1.3.2} If
\[
1\leq p,q\leq\infty \quad \text{and}\quad
\frac{1}{p}-\frac{1}{q}<\frac{\delta N_{\min}(k)}{Q},
\]
then the $\ell^{q,p}$ cohomology in degree $k$ of $G$ does not vanish.
\end{enumerate}
\end{theo}

The non-vanishing statement has a wider scope, see Theorem~\ref{thm:nonzero}. It holds in particular on more general homogeneous groups.

Theorem~\ref{1} is sharp when both $H^{k-1}(\mathfrak{g})$ and $H^{k}(\mathfrak{g})$ are concentrated in a single weight. This happens in all degrees for abelian groups and for Heisenberg groups, for instance. This happens for all Carnot groups in degrees $1$ and $n$: $\delta N_{\max}(1)=\delta N_{\min}(1)=1=\delta N_{\max}(n)=\delta N_{\min}(n)$.

Even when not sharp, the result seems of some value as it relates large scale quasi-isometric analytic invariants of $G$ to its infinitesimal Lie structure. For instance for $k=2$, the weights on $H^2(\mathfrak{g})$ can be interpreted as the depth of the relations defining $G$ with respect to a free Lie group over $\mathfrak{g}_1$, see e.g.~\cite{Rumin_TSG}. In particular, one has always $1\leq \delta N_{\min}(2) \leq \delta N_{\max}(2) \leq r$, where $r$ is the depth of $G$.

%\smallskip

We shall also illustrate the results on the Engel group in Section~\ref{sec:engels-group}, and show in particular that, apart from degrees $1$ and $n$, the natural Carnot homogeneous structure does not give the best range for the non-vanishing result in general.

%\smallskip

Coming back to the general vanishing statement in Theorem~\ref{1}, it can be notice that this bound is reminiscent of former works of the second author on Novikov Shubin invariants~\cite{Rumin_TSG, Rumin10}. Roughly speaking, these numbers can be defined on Galois $\Gamma$-coverings $\widetilde M$ of compact manifolds (or finite simplicial complexes) $M = \widetilde M / \Gamma$, using the notion of von Neumann $\Gamma$-trace of $\Gamma$-invariant operators on $\widetilde M$, see e.g.~\cite[Section~1]{Rumin_TSG} for an introduction. The $k$-th Novikov Shubin number of $M$ measures the asymptotic polynomial decay of the $\Gamma$-trace of spectral projectors $\Pi_{\delta \dd} (]0, \lambda])$ when $\lambda $ goes to $0$ on $k$-forms (or discrete cochains),
\begin{equation*}
\alpha_k(M) = 2 \liminf_{\lambda \rightarrow 0}\bigl(\ln
\Trace_\Gamma (\Pi_{\delta \dd} (\mathopen{]}0, \lambda]) / \ln \lambda \bigr).
\end{equation*}
It turns out that these numbers are related to the $\ell^{q,2}$ cohomology of $\widetilde M$. Namely, following~\cite[Theorem~1.10]{Rumin10}, it holds that the $\ell^{q,2}$ cohomology of $\widetilde M$ vanishes in degree $k$ if
\begin{equation}\label{eq:20}
\frac{1}{2} - \frac{1}{q} \geq \frac{1}{\alpha_{k-1}(M)}.
\end{equation}
Moreover, from~\cite[Theorem~3.13]{Rumin_TSG} one has on rational Carnot groups
\begin{equation*}
\frac{\delta N_{\max}(k)}{Q} \geq \frac{1}{\alpha_{k-1}(G)},
\end{equation*}
if $w_{\min}(k-1)= w_{\max}(k-1)$ in degree $k-1$. Hence, one recovers the vanishing result in Theorem~\ref{1}, for $p=2$, under the assumptions that $H^{k-1}(\mathfrak{g})$ is homogeneous and $G$ is rational, i.e. possesses a cocompact discrete group. However, it should be emphasized that a result using $\Gamma$-dimension like~\eqref{eq:20} relies on (more elementary) Hilbertian techniques, specific to the case $p=2$, and looks difficult to extend this way for other values of $p$, at last on forms when Markovian properties of Laplacian are not available.

\subsection{Method}\label{sec:method}

We now briefly describe the scheme of the proof of Theorem~\ref{1}. The first step is a Leray type lemma which relates the discrete $\ell^{q,p}$ cohomology to some Sobolev $L^{q,p}$ cohomology of differential forms. This is proved here in Theorem~\ref{prop:leray} in the more general setting of manifolds of bounded geometry of some high order. One has to take care of an eventual lack of uniformity in the coverings in order to be able to use local Poincar\'e inequalities.

A feature of this Sobolev $L^{q,p}$ cohomology is that its forms are
\emph{a priori} quite smooth, and need only be integrated into much less regular ones. This is because $\ell^p$-cochains at the discrete level transform into smooth forms made from a smooth partition of unity, while reversely, less regular $L^p$ forms can still be discretized into some $\ell^p$ data.

We then focus on Carnot groups. Although they possess dilatations, the de Rham differential is not homogeneous when seen as a differential operator. Nevertheless, as observed in~\cite{Rumin_CRAS_1999}, it has some type of graded hypoellipticity, that can be used to produce global homotopies $K$. These homotopies are pseudodifferential operators as studied by Folland and Christ--Geller--G{\l}owacki--Polin in~\cite{CGGP,Folland_1975}. They can be thought of as a kind of generalized Riesz potentials, like $\Delta^{-1}\delta$ on $1$-forms, but adapted here to the Carnot homogeneity of the group.

One needs then to translate the graded Sobolev regularity of $K$ into a standard one to get the $L^{q,p}$ Sobolev controls on $\dd$. This is here that the weight gaps of forms arise to control the $(p,q)$ range. Actually, in order to reduce the gap to cohomology weights in $H^*(\mathfrak{g})$ only, we work with a \emph{contracted de Rham complex} $\dd_c$ instead, available in Carnot geometry. It shares the same graded analytic regularity as $\dd$, but uses forms with retracted components over $H^*(\mathfrak{g})$ only, see Section~\ref{sec:dc_complex}. It is worthwhile noting that although the retraction of de Rham complex on $\dd_c$ costs a lot of derivatives, this is harmless here due to the feature of Sobolev $L^{q,p}$ cohomology we mentioned above. Actually, in this low energy large scale problem, loosing regularity is not an issue and the lower derivatives the homotopy $K_c$ controls, the smaller is the $(p,q)$ gap in Sobolev inequality, and the better becomes the $\ell^{q,p}$ vanishing result.

%\smallskip

The non-vanishing result~\eqref{1.3.2} in Theorem~\ref{1} relies on the construction of homogeneous closed differential forms of any order and controlled weights. The contracted de Rham complex is useful to this end too. Such homogeneous forms belong to $L^p$ space with explicit $p$, but can not be integrated in $L^q$ for $q$ too close to $p$, as seen using Poincar\'e duality (construction of compactly supported test forms and integration by parts).


\section{Local Poincar\'e inequality}

In this section, differential forms of degree $-1$, $\Omega^{-1}$, are meant to be constants. The complex is completed with the map $d:\Omega^{-1}\to\Omega^0$ which maps a constant to a constant function with the same value.

\begin{defi}
Let $M$ be a Riemannian manifold. Say that $M$ has $C^\ell$-bounded geometry if injectivity radius is bounded below and curvature together with all its derivatives up to order $\ell$ are uniformly bounded. For $h \leq \ell-1$, let $W^{h,p}$ denote the space of smooth functions $u$ on $M$ which are in $L^p$ as well as all their covariant derivatives up to order $h$. When $p=\infty$, $W^{h,\infty}=C^h\cap L^{\infty}$.
\end{defi}

\begin{rema}\label{whpdecr}
On a $C^\ell$-bounded geometry $n$-manifold, these Sobolev spaces are interlaced. Let $p \leq q$ and $\frac{n}{p} +1 \leq h \leq \ell-1$. Then $W^{h,p}\subset W^{h -1- n/p,q}$.
\end{rema}

Indeed, on a ball $B$ of size smaller than the injectivity radius, usual Sobolev embedding holds, $W^{h,p}(B)\subset L^q(B)$ provided $\frac{1}{p}-\frac{1}{q}\leq \frac{h}{n}$, an inequality which is automatically satisfied if $h\geq \frac{n}{p}+1$. Pick a covering $B_i$ of $M$ by such balls with bounded multiplicity. Let $u_i=\|f\|_{W^{h,p}(B)}$ and $v_i=\|f\|_{W^{h-1 -n/p,q}(B)}$. Then $v_i\leq C\,u_i$. Furthermore
\[
\|f\|_{W^{h-1 -n/p,q}}\leq C\,\|(v_i)\|_{\ell^q}\leq C\,\|(v_i)\|_{\ell^p}\leq C'
\|(u_i)\|_{\ell^p}\leq C''\,\|f\|_{W^{h,p}}.
\]

\begin{rema}
According to Bemelmans--Min-Oo--Ruh~\cite{B-MinOo-Ruh}, for any fixed $h$, any complete Riemannian metric with bounded curvature can be approximated by an other one with all derivatives of curvature up to order $h$ uniformly bounded.
\end{rema}

Therefore assuming bounded geometry up to a high order is not a restriction for our overall purposes. The main point is that curvature and injectivity radius be bounded. Nevertheless, it helps in technical steps like the following Proposition.

\begin{prop}\label{local}
Let $M$ be a Riemannian manifold with sectional curvature bounded by $K$, as well as all covariant derivatives of curvature up to order $\ell$. Let $R<\frac{\pi}{2\sqrt{K}}$. Assume that $M$ has a positive injectivity radius, larger than $2R$. Let $y\in M$. Let $U_j$ be balls in $M$ containing $y$ of radii $\leq R$, and $U=\bigcap_j U_j$. The Cartan homotopy is an operator $P$ on differential forms on $U$ which satisfies $1=P\dd+\dd P$ and maps $C^h\Omega^{k}(U)$ to $C^h\Omega^{k-1}(U)$ for all $h\leq \ell -1$ and all $k=0,\dots,n$, with norm depending on $K$, $h$ and $R$ only.

Assume further that $U$ contains $B(y,r)$ for some $r>0$. Then $P$ is bounded from $W^{h,p}\Omega^{k}(U)$ to $W^{h -n-1,q}\Omega^{k-1}(U)$, provided $p\geq 1$, $q\geq 1$, $n+1 \leq h \leq \ell-1$. Its norm depends on $K$, $\ell$, $n$, $R$ and $r$ only.
\end{prop}

\begin{proof}
The assumptions on $R$ guarantee that minimizing geodesics between points at distance $<R$ are unique and that all balls of radii $\leq R$ are geodesically convex. For $x,y\in U$, let $\gamma_{x,y}$ denote the unique minimizing geodesic from $x$ to $y$, parametrized on $[0,1]$ with constant speed $\dd(x,y)$. Fix $y\in U$. Consider the vectorfield $\xi_y$ defined as follows,
\[
\xi_{y}(x)=\gamma'_{x,y}(0).
\]
It is smooth, since in normal coordinates centered at $y$, $\xi_{y}$ is the radial vectorfield $\xi_y(u)=-u$. Let $\phi_{y,t}$ denote the diffeomorphism semi-group generated by $\xi_y$. For $t\in\R_+$, $\phi_{y,t}$ maps a point $x$ to $\gamma_{y,x}(e^{-t}\dd(y,x))$.

Let $k\geq 1$. Following H. Cartan, define an operator $P_y$ on $k$ forms $\omega$ by
\[
P_y(\omega)=-\int_{0}^{+\infty}\phi_{y,t}^{*}\iota_{\xi_{y}}\omega\,\dt.
\]
Then, on $k$-forms,
\begin{align*}
\dd P_y+P_y\dd &=- \int_{0}^{+\infty}\phi_{y,t}^{*} (\dd\iota_{\xi_{y}} +
\iota_{\xi_{y}}\dd)\,\dt\\
&=-\int_{0}^{+\infty}\phi_{y,t}^{*}\mathcal{L}_{\xi_{y}}\,\dt\\
&=-\phi_{y,+\infty}^{*}+\phi_{y,0}^{*}=1,
\end{align*}
if $k\geq 1$, since $\phi_{y,0}$ is the identity and $\phi_{y,+\infty}$ is the constant map to $y$. On $0$-forms, define $P_y(\omega)=\omega(y)$. Then $\dd P_y=\omega(y)$ viewed as a (constant) function on $U$, whereas $P_y \dd\omega=\omega-\omega(y)$, hence $\dd P_y+P_y \dd=1$ also on $0$-forms.

In normal coordinates with origin at $y$, $P_y$ has a simple expression
\begin{align*}
P_y\omega(x)&=\int_{0}^{+\infty}e^{-kt}\omega_{e^{-t}x}(x,\dots)\,\dt\\
&=\int_{0}^{1}s^{k-1}\omega_{sx}(x,\dots)\,\ds.
\end{align*}
It shows that $P_y$, read in normal coordinates, is bounded on $C^h$ for all $h$.

The domain of exponential coordinates, $V=\exp_y^{-1}(U)$, is convex. If it contains a ball of radius $r$, there is a bi-Lipschitz homeomorphism of the unit ball to $V$ with Lipschitz constants depending on $R$ and $r$ only, hence Sobolev embeddings $W^{1,p}\subset L^q$ for $\frac{1}{p}-\frac{1}{q}\leq\frac{1}{n}$ with uniform constants (if $q=\infty$, $p>n$ is required and $L^q$ is replaced with $C^0$). If $p\geq 1$, this implies that $W^{n+1,p}\subset C^0$, hence $W^{h,p}\subset C^{h-n-1}$. Obviously, $C^{h-n-1}\subset W^{h-n-1,q}$.

Since curvature and its derivatives are bounded up to order $\ell$, the Riemannian exponential map and its inverse are $C^{\ell-1}$-bounded, hence $P_y$ is bounded on $C^h$ for $h \leq \ell-1$. If furthermore $h \geq n+1$, the embeddings
\[
W^{h,p}\subset C^{h-n-1}\subset W^{h-n-1,q}
\]
hold on $U$ with bounds depending on $K$, $R$ and $r$ only, hence $P$ maps $W^{h,p}$ to $W^{h-n-1,q}$, with uniform bounds.
\end{proof}


\section{\texorpdfstring{$\ell^{q,p}$}{ell(q,p)} cohomology and Sobolev \texorpdfstring{$L^{q,p}$}{L(q,p)} cohomology}\label{sec:leray}

\begin{defi}
Let $M$ be a Riemannian manifold. Let $h, h'\in\N$. The
\emph{Sobolev $L^{q,p}$ cohomology} is
\[
L_{h',h}^{q,p}H^{\cdot}(M)=\{\text{closed forms in }W^{h,p}\} / \dd(\{\text{forms in }W^{h',q}\}) \cap W^{h,p}.
\]
\end{defi}

\begin{rema}\label{rem:decreasing}
On a bounded geometry $n$-manifold, this is nonincreasing in $q$ in the following sense. Let $q'\geq q$ and $h'\geq n/q + 1$. Then $L_{h',h}^{q,p}H^{k}(M)$ surjects onto $L_{h'-1 - n/q,h}^{q',p}H^{k}(M)$, see Remark~\ref{whpdecr}.
\end{rema}

\begin{theo}\label{prop:leray}
Let $1\leq p\leq q\leq \infty$. Let $M$ be a Riemannian manifold of $C^\ell$-bounded geometry, with $\ell > n(n+1)$. There exists a simplicial complex $T$, admitting a quasiisometric homotopy equivalence to $M$, with the following property. For every integers $h,h'$ such that $n(n+1) \leq h,h'\leq \ell -1$, there exists an isomorphism between the $\ell^{q,p}$ cohomology of $T$ and the Sobolev $L^{q,p}$ cohomology $L_{h',h}^{q,p}H^{\cdot}(M)$.
\end{theo}

The proof is a careful inspection of Leray's acyclic covering theorem.

First construct a simplicial complex $T$ quasiisometric to $M$. Up to rescaling, one can assume that sectional curvature is $\leq 1/n^2$ and injectivity radius is $>2n$. Pick a maximal $1/2$-separated subset $\{x_i\}$ of $M$. Let $B_i$ be the covering by closed unit balls centered on this set. Let $T$ denote the nerve of this covering. Let $U_i=B(x_i,3)$. Note that if $x_{i_0},\dots,x_{i_j}$ span a $j$-simplex of $T$, then the intersection
\[
U_{{i_0}\dots{i_j}} := \bigcap_{m=0}^{j}U_{i_m}
\]
is contained in a ball of radius 3 and contains a concentric ball of radius 1.

Pick once and for all a smooth cut-off function with support in $[-1,1]$, compose it with distance to points $x_i$ and convert the obtained collection of functions into a partition of unity $\chi_i$ by dividing by the sum.

Define a bicomplex $C^{j,k}=$ skew-symmetric maps associating to $j+1$-tuples $(i_0,\dots, i_j)$ differential $k$-forms on $j+1$-fold intersections $U_{i_0}\cap\dots\cap U_{i_j}$. It is convenient to extend the notation to
\[
C^{-1,k} =\Omega^k(M),\quad C^{j,-1}=C^j(T),\quad C^{j,\cdot}=C^{\cdot,j}=0 \text{ if } j<-1.
\]

The two commuting complexes are $\dd:C^{j,k}\to C^{j,k+1}$ and the simplicial coboundary $\delta: C^{j,k}\to C^{j+1,k}$ defined by
\[
\delta(\phi)_{i_0 \dots i_{j+1}}= \phi_{i_0 \dots i_{j}}-\phi_{i_0
\dots i_{j-1}j_{j+1}} +\dots+(-1)^{j+1}\phi_{i_1 \dots i_{j+1}},
\]
restricted to $U_{i_0}\cap\dots\cap U_{i_{j+1}}$. By convention, $\dd:C^{j,-1}\to C^{j,0}$ maps scalar $j$-cochains to skew-symmetric maps to functions on intersections which are constant. Also, $\delta:C^{-1,k}\to C^{0,k}$ maps a globally defined differential form to the collection of its restrictions to open sets $U_i$. We define the differential to be zero in all other cases.
\[
\xymatrix %@!=1pc
{ & C^{-1,0}\ar[d]_{\delta} & C^{-1,1}\ar[d]_{\delta} & C^{-1,2}\ar[d]_{\delta} & C^{-1,3}\ar[d]_{\delta} &\dots \\ C^{0,-1} \ar[r]^{d} & C^{0,0}
\ar[r]^{d}\ar[d]_{\delta} & C^{0,1}\ar[r]^{d}\ar[d]_{\delta} & C^{0,2}\ar[r]^{d}\ar[d]_{\delta} & C^{0,3}\ar[r]^{d}\ar[d]_{\delta} &\dots
\\
C^{1,-1} \ar[r]^{d} & C^{1,0} \ar[r]^{d}\ar[d]_{\delta} & C^{1,1}
\ar[r]^{d}\ar[d]_{\delta} & C^{1,2} \ar[r]^{d}\ar[d]_{\delta} & C^{1,3} \ar[r]^{d}\ar[d]_{\delta} &\dots
\\
C^{2,-1} \ar[r]^{d} & C^{2,0} \ar[r]^{d}\ar[d]_{\delta} & C^{2,1}
\ar[r]^{d}\ar[d]_{\delta} & C^{2,2} \ar[r]^{d}\ar[d]_{\delta} & C^{2,3} \ar[r]^{d}\ar[d]_{\delta} &\dots
\\
\dots &\dots &\dots &\dots &\dots &\dots }
\]

The coboundary $\delta$ is inverted by the operator
\[
\epsilon:C^{j,k}\to C^{j-1,k}
\]
defined by
\[
\epsilon(\phi)_{i_0 \dots i_{j-1}}=\sum_{m}\chi_m \phi_{m i_0 \dots i_{j-1}}.
\]
By an inverse, we mean that $\delta\epsilon+\epsilon\delta=1$. This identity persists in all nonnegative bidegrees provided $\epsilon : C^{0,k} \to C^{-1,k}$ is defined by $\epsilon (\phi) = \sum_m \chi_m \phi_m$ and $\epsilon=0$ on $C^{j,k}$, $j<0$.
\[
\xymatrix %@!=1pc
{ & C^{-1,0} & C^{-1,1} & \textcolor{red}{C^{-1,2}} & C^{-1,3} &\dots
\\
C^{0,-1} & C^{0,0} \ar[l]^{P}\ar[u]^{\epsilon} & C^{0,1}\ar[l]^{P}\ar[u]^{\epsilon} & C^{0,2} & C^{0,3}\ar[l]^{P}\ar[u]^{\epsilon} &\dots\ar[l]^{P}
\\
C^{1,-1} & C^{1,0} \ar[l]^{P}\ar[u]^{\epsilon} & C^{1,1} & C^{1,2} \ar[l]^{P}\ar[u]^{\epsilon} & C^{1,3} \ar[l]^{P}\ar[u]^{\epsilon} &\dots\ar[l]^{P}
\\
\textcolor{red}{C^{2,-1}} \ar@{=>>} '[r]^-{\textcolor{red}{d}} '[ru]^{\textcolor{red}{\epsilon}} '[rru]^{\textcolor{red}{d}} '[rruu]^{\textcolor{red}{\epsilon}} '[rrruu]^{\textcolor{red}{d}} '[rrruuu]^{\textcolor{red}{\epsilon}}& C^{2,0} & C^{2,1}
\ar[l]^{P}\ar[u]^{\epsilon} & C^{2,2} \ar[l]^{P}\ar[u]^{\epsilon} & C^{2,3} \ar[l]^{P}\ar[u]^{\epsilon} &\dots\ar[l]^{P}
\\
\dots &\dots\ar[u]^{\epsilon} &\dots\ar[u]^{\epsilon} &\dots\ar[u]^{\epsilon} &\dots\ar[u]^{\epsilon} &\dots }
\]

Consider the maps
\[
\Phi_j = (\epsilon \dd)^{j+1} : C^j(T)=C^{j,-1} \to C^{-1,j}=\Omega^j(M).
\]
($\Phi_2$ is featured on the diagram). By definition, given a cochain $\kappa$, $\Phi_j(\kappa)$ is a local linear expression of the constants defining $\kappa$, multiplied by polynomials of the $\chi_i$ and their differential. Therefore $\Phi_j(\kappa)$ is $C^\infty$ and belongs to $W^{h,q}(\Omega^j(M))$ for all $q \geq p\geq 1$ if $\kappa \in \ell^p(C^j(T))$.

Since
\[
(\epsilon \dd)\delta=\epsilon\delta \dd=(1-\delta\epsilon)\dd=\dd-\delta\epsilon \dd,
\]
for $j\geq 0$, on $C^{j,-1}$,
\begin{align*}
(\epsilon \dd)^{j+2}\delta &=(\epsilon \dd)^{j+1} (\dd-\delta\epsilon \dd)= -(\epsilon \dd)^{j+1} \delta(\epsilon \dd) \\
&=\dots\\
&=(-1)^{j+1}(\epsilon \dd)\delta(\epsilon \dd)^{j+1}=(-1)^{j+1}(\dd-\delta\epsilon \dd) (\epsilon \dd)^{j+1}\\
&=(-1)^{j+1}\dd(\epsilon \dd)^{j+1}-(-1)^{j+1}\delta(\epsilon \dd)^{j+2}\\
&=(-1)^{j+1}\dd(\epsilon \dd)^{j+1}.
\end{align*}
Indeed, $(\epsilon \dd)^{j+2}(C^{j,-1})\subset C^{-2,j+1}=\{0\}$. In other words,
\[
\Phi_{j+1} \circ \delta=(-1)^{j+1}\dd\circ\Phi_j.
\]

%\medskip

We now proceed in the opposite direction to produce a cohomological inverse of $\Phi$. Let us mark each intersection $U_{i_0 \dots i_j}:=U_{i_0}\cap\dots\cap U_{i_j}$ with the point $y=x_{i_0}$. Proposition~\ref{local} provides us with an operator $P_{i_0\dots i_j}$ on usual $k$-forms on $U_{i_0 \dots i_j}$. Putting them together yields an operator $P:C^{j,k}\to C^{j,k-1}$ such that $1=\dd P+P\dd$. Furthermore, $P$ is bounded from $W^{\cdot,p}$ to $W^{\cdot-n-1,q}$.

Exchanging the formal roles of $(\delta, \epsilon)$ and $(\dd,P)$, we define
\[
\Psi_k = (P\delta)^{k+1} :\Omega^k(M)=C^{-1,k}\to C^{k,-1}=C^k(T),
\]
As above, one checks easily that $\Psi_{k+1} \circ \dd = (-1)^{k+1}\delta \circ \Psi_k$.

Observe that the maps $\Psi_j \circ \Phi_j$, $j=0,1,\dots$, put together form a morphism of the complex $C^{\cdot,-1}=C^\cdot(T)$ (i.e. they commute with $\delta$). We next show that it is homotopic to the identity. Let us prove by induction on $i$ that, on $C^{\cdot,-1}$,
\begin{equation}\label{eq:7}
(P \delta)^i (\epsilon \dd)^i = 1 - R_i \delta - \delta R_{i-1}\,,
\end{equation}
with $R_0= 0$ and $R_i = \sum_{k=0}^{i-1} (-1)^k (P \delta)^k P (\epsilon \dd)^{k+1}$. This implies the result for $\Psi_j \circ \Phi_j$.
\begin{proof}
For $i=1$, one has on $C^{\cdot,-1}$
\begin{align*}
(P \delta) (\epsilon \dd) & = P(1 - \epsilon \delta)\dd = P\dd - P\epsilon \dd \delta \\
& = 1 - (P\epsilon \dd) \delta,
\end{align*}
since $P\dd = 1$ on $C^{\cdot,-1}$. Assuming~\eqref{eq:7} for $i$, one writes
\begin{align}
(P\delta)^{i+1} (\epsilon \dd)^{i+1} & = (P \delta)^i P\delta
\epsilon \dd (\epsilon \dd)^i = (P \delta)^i P(1 - \epsilon
\delta) \dd (\epsilon \dd)^i \nonumber \\
& = (P
\delta)^i (1 - \dd P - P \epsilon \delta \dd) (\epsilon \dd)^i
\nonumber \\
& = (P \delta)^i (\epsilon \dd)^i - (P \delta)^i \dd P (\epsilon \dd)^i - (P \delta)^i P (\epsilon \dd) \delta (\epsilon \dd)^i.\label{eq:8}
\end{align}

About the second term in~\eqref{eq:8}, one finds that $\delta(P\delta) \dd = - \delta \dd (P\delta)$, hence by induction, one can push the isolated $\dd$ term to the left,
\[
(P \delta)^i d = (-1)^{i-1} P\delta \dd (P\delta)^{i-1} = (-1)^{i-1} P\dd \delta (P\delta)^{i-1} = (-1)^{i-1}\delta (P\delta)^{i-1},
\]
when the image lies within $C^{\cdot,-1}$, as it does in~\eqref{eq:8}.

For the third term in~\eqref{eq:8}, one sees that $(\epsilon \dd) \delta (\epsilon \dd) = (\epsilon \dd)(1 - \epsilon \delta) \dd = - (\epsilon \dd)^2 \delta$, so that one can push the isolated $\delta$ term to the right,
\[
(P \delta)^i P (\epsilon \dd) \delta (\epsilon \dd)^i = (-1)^i (P \delta)^i P (\epsilon \dd)^{i+1} \delta \,.
\]
Gathering in~\eqref{eq:8} gives
\[
(P\delta)^{i+1} (\epsilon \dd)^{i+1} = 1 - (R_i + (-1)^i (P\delta)^i P (\epsilon \dd)^{i+1}) \delta - \delta (R_{i-1} + (-1)^{i-1} (P\delta)^{i-1} P (\epsilon \dd)^i),
\]
that proves~\eqref{eq:7}.
\end{proof}

Similarly, $\Phi\circ\Psi$ is homotopic to the identity on the complex $(C^{-1,\cdot},\dd)$, $\Phi\circ\Psi=1-\dd R'-R'\dd$.

Finally, let us examine how Sobolev norms behave under the class of endomorphisms we are using. Maps from cochains to differential forms, i.e. $\epsilon$, $\Phi$ are bounded from $\ell^p$ to $W^{\ell-1,p}$. Maps from differential forms, i.e. $P$, $\Psi$, loose derivatives (but this is harmless since the final outputs are scalar cochains) but are bounded from $W^{h,p}$ to $\ell^p$, $h\leq \ell-1$. One merely needs $h$ large enough to be able to apply $P$ $n$ times, whence the assumption $h\geq n(n+1)$. Maps from cochains to cochains, e.g. $R$, are bounded on $\ell^p$, maps from differential forms to differential forms, e.g. $R'$, are bounded from $W^{h',p}$ to $W^{h,p}$ for every $h'\geq n(n+1)$ such that $h'\leq \ell-1$.

If $q\geq p$, the $\ell^q$-norm is controlled by the $\ell^p$-norm, hence $R$ is bounded from $\ell^p$ to $\ell^q$. It is also true that $R'$ is bounded from $W^{h'-1,p}$ to $W^{h-1,q}$. Indeed, it is made of bricks which map differential forms to cochains or cochains to differential forms, so no loss on derivatives affects the final differentiability. For the same reason, one can gain local integrability from $L^p_{loc}$ to $L^q_{loc}$ without restriction on $p$ and $q$ but $p$, $q\geq 1$.

It follows that $\Phi$ and $\Psi$ induce isomorphisms between the $\ell^{q,p}$ cohomology of $T$ and the Sobolev $L^{q,p}$ cohomology.


\section{De Rham complex and graduation on Carnot groups}\label{sec:dc_complex}

From now on, we will work on Carnot Lie groups. These are nilpotent Lie groups $G$ such that their Lie algebra $\mathfrak{g}$ splits into a direct sum
\[
\mathfrak{g} = \mathfrak{g}_1 \oplus \mathfrak{g}_2 \oplus \dots
\oplus\mathfrak{g}_r \quad \mathrm{satisfying} \quad [\mathfrak{g}_1,
\mathfrak{g}_i] = \mathfrak{g}_{i+1} \ \mathrm{for} \ 1\leq i \leq r-1.
\]
The weight $w = i$ on $\mathfrak{g}_i$ induces a family of dilations $\delta_t = t^w$ on $\mathfrak{g} \simeq G$.

In turn, the tangent bundle $TG$ splits into left invariant sub-bundles $H_1 \oplus \dots \oplus H_r$ with $H_i = \mathfrak{g}_i$ at the origin. Finally, differential forms decompose through their weight $\Omega^k G = \oplus_w \Omega^{k,w} G$ with $\Omega^{k_1} H^*_{w_1}\wedge \dots \wedge \Omega^{k_i}H^*_{w_i}$ of weight $w = k_1 w_1 + \dots + k_i w_i$. De Rham differential $\dd$ itself splits into
\[
\dd = \dd_0 + \dd_1 + \dots + \dd_r,
\]
with $\dd_i$ increasing weight by $i$. Indeed, this is clear on functions where $\dd_0 = 0$ and $\dd_i = d $ along $H_i$. This extends to forms, using $\dd(f\alpha) = \dd f \wedge \alpha + f \dalpha$ and observing that for left invariant forms $\alpha$ and left invariant vectors $X_i$, Cartan's formula reads
\begin{align*}
\dalpha (X_1,\dots, X_{k+1}) & = \sum_{1\leq i < j\leq k+1} (-1)^{i+j}
\alpha ([X_i,X_j], \dots, \widehat{X}_{i,j},
\dots X_{k+1}) \\
& = \dd_0 \alpha (X_1,\dots, X_{k+1}).
\end{align*}
Hence $\dd= \dd_0$ is a weight preserving algebraic (zero order) operator on invariant forms, and $\ker \dd_0 / \im \dd_0 = H^*(\mathfrak{g})$ is the Lie algebra cohomology of $G$. Note also that from these formulas, $\dd_i$ is a homogeneous differential operator of degree~$i$ and increases the weight by $i$. Like $\dd$, it is homogeneous of degree~$0$ through $h^*_\lambda$, since $\dd h^*_\lambda = h^*_\lambda \dd$.

%\medskip

This algebraic $\dd_0$ allows to split and contract de Rham complex on a smaller subcomplex, as we now briefly describe. This was shown in~\cite{Rumin_CRAS_1999} and~\cite{Rumin_TSG} in the more general setting of Carnot--Caratheodory manifolds. More details may be found there.

Pick an invariant metric so that the $\mathfrak{g}_i$ are orthogonal to each others, and let $\delta_0 = \dd_0^*$ and $\dd_0^{-1}$ be the partial inverse of $\dd_0$ such that $\ker \dd_0^{-1} = \ker \delta_0$, $\dd_0^{-1} \dd_0 = \Pi_{\ker \delta_0}$ and $\dd_0 \dd_0^{-1} = \Pi_{\im \dd_0}$.

Let $E_0 = \ker \dd_0 \cap \ker \delta_0 \simeq \Omega^* H^*(\mathfrak{g})$. Iterating the homotopy $r = 1 - \dd_0^{-1}\dd - \dd \dd_0 ^{-1}$ one can show the following results, stated here in the particular case of Carnot groups.

\let\oldpointir\pointir
\renewcommand{\pointir}{}
\begin{theo}[{\cite[Theorem~1]{Rumin_CRAS_1999}}]\label{thm:dc-complex}
\ %\ \\*[-2.6em]
\begin{enumerate}
\item The de Rham complex on $G$ splits as the direct sum of two sub-complexes $E \oplus F$, where $E = \ker \dd_0^{-1} \cap \ker \dd \dd_0^{-1}$ and $F= \im \dd_0^{-1} + \im (\dd\dd_0^{-1}) $.
\item The retractions $r^k$ stabilize to $\Pi_E$ the projection on $E$ along $F$. $\Pi_E$ is a homotopy equivalence of the form $\Pi_E = 1 - R\dd - \dd R$ where $R$ is a differential operator.
\item One has $\Pi_E \Pi_{E_0} \Pi_E = \Pi_E$ and $\Pi_{E_0} \Pi_E \Pi_{E_0} = \Pi_{E_0}$ so that the complex $(E,\dd)$ is conjugated through $\Pi_{E_0}$ to the complex $(E_0, \dd_c)$ with $\dd_c = \Pi_{E_0} \dd \Pi_E \Pi_{E_0}$.
\end{enumerate}
\end{theo}
\let\pointir\oldpointir

This shows in particular that the de Rham complex, $(E,\dd)$ and $(E_0, \dd_c)$ are homotopically equivalent complexes on smooth forms. For convenience in the sequel, we will refer to $(E_0,\dd_c)$ as the \emph{contracted de Rham complex} (also known as Rumin complex) and sections of $E_0$ as \emph{contracted forms}, since they have a restricted set of components with respect to usual ones.

We shall now describe its analytical properties we will use.


\section{Inverting \texorpdfstring{$\dd_c$}{dc} and \texorpdfstring{$\dd$}{d} on \texorpdfstring{$G$}{G}}\label{sec:inverting-d-d_c}

De Rham and contracted de Rham complexes are not homogeneous as differential operators, but are indeed invariant under the dilations $\delta_t$ taking into account the weight of forms. This leads to a notion of sub-ellipticity in a graded sense, called C-C ellipticity in~\cite{Rumin_CRAS_1999, Rumin_TSG}, that we now describe.

Let $\nabla = \dd_1$ the differential along $H = H_1$. Extend it on all forms using $\nabla (f \alpha) = (\nabla f) \alpha$ for left invariant forms $\alpha$ on $G$. Kohn's Laplacian $\Delta_H = \nabla^* \nabla$ is hypoelliptic since $H$ is bracket generating on Carnot groups, and positive self-adjoint on $L^2$. Let then
\[
|\nabla| = \Delta_H^{1/2}
\]
denotes its square root. Following~\cite[Section~3]{Folland_1975} or~\cite{CGGP}, it is a homogeneous first order pseudodifferential operator on $G$ in the sense that its distributional kernel, acting by group convolution, is homogeneous and smooth away from the origin. It possesses an inverse $|\nabla|^{-1}$, which is also a homogeneous pseudodifferential operator of order $-1$ in this calculus. Actually according to~\cite[Theorem~3.15]{Folland_1975}, it belongs to a whole analytic family of pseudodifferential operators $|\nabla|^\alpha$ of order $\alpha \in \mathbb{C}$. Note that kernels of these homogeneous pseudodifferential operators may contain logarithmic terms, when the order is an integer $\leq -Q$. We refer to~\cite{CGGP} and~\cite[Section~1]{Folland_1975} for more details and properties of this calculus.

%\smallskip
A particularly useful test function space for these operators is given by the space of Schwartz functions all of whose polynomial moments vanish,
\begin{equation}\label{eq:9}
\calS_0=\{f\in\calS\,;\,\langle f,P\rangle=0 \text{ for every polynomial }P\},
\end{equation}
where $\mathcal{S}$ denotes the Schwartz space of $G$ and $\langle f,P\rangle=\int_Gf(x)P(x)\,\dx$. Unlike more usual test functions spaces as $C^\infty_c$ or $\mathcal{S}$, this space $\calS_0$ is stable under the action of pseudodifferential operators of any order in the calculus, see~\cite[Proposition~2.2]{CGGP}, so that they can be composed on it. In particular by~\cite[Theorem~3.15]{Folland_1975}, for every $\alpha$, $\beta \in \mathbb{C}$,
\begin{equation}\label{eq:10}
|\nabla|^\alpha |\nabla|^{\beta} = |\nabla|^{\alpha+\beta} \
\mathrm{on} \ \calS_0 \,.
\end{equation}
We shall prove in Proposition~\ref{prop:density-s_0} that $\calS_0$ is dense in all Sobolev spaces $W^{h,p}$ with $h \in \N$ and $1< p< +\infty$, but we shall work mainly in $\calS_0$ in this section.

%\medskip

Now let $N=w$ on forms of weight $w$. Consider the operator $|\nabla|^N$, preserving the degree and weight of forms, and acting componentwise on $\calS_0$ in a left-invariant frame.

From the previous discussion, $\dd^{\nabla} = |\nabla|^{-N} \dd |\nabla|^N$ and $\dd_c^{\nabla} = |\nabla|^{-N} \dd_c |\nabla|^N$ are both homogeneous pseudodifferential operators of (differential) order $0$. Indeed, as observed in Section~\ref{sec:dc_complex}, $\dd$ splits into $\dd=\sum_i \dd_i$ where $\dd_i$ is a differential operator of horizontal order $i$ which increases weight by $i$. On forms of weight $w$, $|\nabla|^N$ has differential order $w$, $\dd_i |\nabla|^N$ has order $w+i$ and maps to forms of weight $w+i$, on which $ |\nabla|^{-N}$ has order $-(w+i)$, hence $|\nabla|^{-N} \dd |\nabla|^N$ has order $0$. The same argument applies to $\dd_c$.

Viewed in this Sobolev scale, these complexes become invertible in the pseudodifferential calculus. Let
\[
\Delta^{\nabla } = \dd^{\nabla} (\dd^{\nabla})^* + (\dd^{\nabla})^* \dd^\nabla
\quad \mathrm{and} \quad \Delta_c^{\nabla} = \dd_c^{\nabla} (\dd_c^{\nabla})^* + (\dd_c^{\nabla})^* \dd_c^\nabla,
\]
not to be confused with the non homogeneous $\dd_c^{\nabla} (\dd_c^*)^{\nabla} + (\dd_c^*)^{\nabla} \dd_c^{\nabla} = |\nabla|^{-N} (\dd_c \dd_c^* + \dd_c^* \dd_c) |\nabla|^N$.

\begin{theo}[{\cite[Theorem~3]{Rumin_CRAS_1999}, \cite[Theorem~5.2]{Rumin_TSG}}]\label{thm:left_inverse}
The Laplacians $\Delta^{\nabla} $ and $\Delta_c^{\nabla } $ have left inverses $Q^{\nabla}$ and $Q_c^{\nabla}$, which are zero order homogeneous pseudodifferential operators.
\end{theo}

By~\cite[Theorem~6.2]{CGGP}, this amounts to show that these Laplacians satisfy Rockland's injectivity criterion. This means that their symbols are injective on smooth vectors of any non trivial irreducible unitary representation of $G$.
%\medskip

This leads to a global homotopy for $\dd_c$ on $G$. Indeed following~\cite[Proposition~1.9]{Folland_1975}, homogeneous pseudodifferential operators of order zero on $G$, such as $\dd_c^{\nabla}$, $\Delta_c^\nabla$ and $Q_c^{\nabla}$, are bounded on all $L^p(G)$ spaces for $1< p< \infty$. Therefore the positive self-adjoint $\Delta_c^{\nabla }$ on $L^2(G)$ is bounded from below since $Q_c^{\nabla} \Delta_c^{\nabla} = 1$. Hence, it is invertible in $L^2(G)$ and $Q_c^{\nabla} = (\Delta_c^{\nabla})^{-1}$ is the inverse of $\Delta_c^{\nabla}$.

Since $(\Delta_c^{\nabla})^{-1}$ commute with $\dd_c^{\nabla}$, the zero order homogeneous pseudodifferential operator $K_c^\nabla = (\dd_c^{\nabla})^* (\Delta_c^{\nabla})^{-1} $ is a global homotopy,
\[
1 = \dd_c^\nabla K_c^\nabla + K_c^\nabla \dd_c^\nabla.
\]
Let us set
\[
K_c = |\nabla|^{N}K_c^\nabla|\nabla|^{-N}
\]
on $\calS_0$, in order that
\[
1 = \dd_c K_c + K_c \dd_c.
\]
Since $K_c^\nabla$ is bounded from $L^p$ to $L^p$ for $1<p< \infty$, $K_c$ is bounded on $\mathcal{S}_0$ endowed with the graded Sobolev norm
\[
\|\alpha\|_{|\nabla|,N,p}:=\||\nabla|^{-N}\alpha\|_p \,.
\]
Actually $K_c$ stays bounded for the whole Sobolev scale with shifted weights $N - m$.

\begin{prop}\label{prop:K_c_scale}
For every constant $m\in\R$ and $1< p< +\infty$, the homotopy $K_c$ is also bounded on $\calS_0$ with respect to the norm $\| \ \|_{|\nabla|,N-m,p}$.
\end{prop}

\begin{proof}
Indeed for $\alpha\in \mathcal{S}_0$,
\begin{align*}
\|K_c \alpha\|_{|\nabla|,N-m,p} &= \||\nabla|^{m-N}K_c\alpha\|_{p}\\
&= \||\nabla|^m | \nabla|^{-N}K_c| \nabla|^{N} |\nabla|^{-m} |\nabla|^{m-N}\alpha\|_{p} \quad \text{by}~\eqref{eq:10},\\
&=\||\nabla|^m K_c^\nabla|\nabla|^{- m}|
\nabla|^{m-N}\alpha\|_{p} \\
&\leq C\, \|| \nabla|^{m-N}\alpha\|_{p} =C\, \|\alpha\|_{|\nabla|,N-m,p} \,,
\end{align*}
since $|\nabla|^m K_c^\nabla|\nabla|^{- m}$ is pseudodifferential of order $0$ and homogeneous, hence bounded on $L^p$ if $1< p< \infty$.
\end{proof}

\begin{rema}\label{rem:1}
Note that using Theorem~\ref{thm:left_inverse}, one can also produce a homotopy in the same way for the full de Rham complex itself. But as we will see in Section~\ref{sec:remarks-weight-gaps}, it leads to a weaker vanishing theorem for $\ell^{q,p}$ cohomology.
\end{rema}


\section{Relating the \texorpdfstring{$|\nabla|$-}{|nabla|-}graded to standard Sobolev norms}

The next step is to compare the $\|\ \|_{|\nabla|, N,p}$ norms, depending on the weight of forms, to usual Sobolev norms of positive order.

Fix a basis $X_i$ of $\mathfrak{g}_1$, viewed as left invariant vectorfields on $G$. Let $\nabla$ denote the horizontal gradient $\nabla f=(X_1 f,\dots,X_{n_1} f)$. Acting componentwise on tuples of functions allows to iterate it. One also extends it on differential forms using their components in a left invariant basis. For $1 \leq p \leq \infty$ and $h \in \N$, define the Sobolev $W_c^{h,p}$ norm
\[
\|\alpha\|_{W_c^{h,p}} = \sum_{k=0}^{h}\|\nabla^k\alpha\|_p \,.
\]

According to Folland~\cite[Theorem~4.10, Corollary~4.13]{Folland_1975}, these norms are equivalent to Sobolev norms defined using $|\nabla|$, provided $1<p<+\infty$. Namely one has then
\begin{equation}\label{eq:11}
\|\ \|_{W_c^{h,p}} \asymp\sum_{k=0}^{h}\||\nabla|^k \ \|_p \,.
\end{equation}

We shall now compare these norms to the graded ones we introduced in the previous section.

\begin{prop}\label{prop:Sobolev}
Let $1< p< \infty$ be fixed, and $a, b, h \in \N$ be such that $a\leq b \leq a+h$. Let $\Omega_{[a,b]}$ denote the space of differential forms whose components have weights $a\leq w \leq b$.
\begin{enumerate}%[label=(\roman*)]
\item \label{6.1.1} It holds on $\Omega_{[a,b]}$ that
\[
\sum_{m=b}^{a+h} \|\ \|_{|\nabla|, N-m, p} \leq
\sum_{k=0}^{h}\||\nabla|^k \ \|_p \leq
\sum_{m=a}^{b+h} \|\ \|_{|\nabla|, N-m, p}\,.
\]

\item \label{6.1.2} Let $\mu \in \N$, $\mu<Q$, and $1< p< q< \infty$ satisfy $\frac{1}{p} - \frac{1}{q} = \frac{\mu}{Q}$. Then for some $C>0$ it holds on $\Omega_{[a,b]} \cap \mathcal{S}$ that
\[
C \| \ \|_{W^{h,q}_c} \leq \sum_{m=a+ \mu}^{b+\mu + h} \| \
\|_{|\nabla|, N-m,p} \,.
\]
\end{enumerate}
\end{prop}

\begin{proofc}
\eqref{6.1.1}\pointir By definition $a\leq N= w \leq b$ on $\Omega_{[a,b]}$, hence
\[
\sum_{m=b}^{a+h} \|\ \|_{|\nabla|, N-m, p} = \sum_{m=b}^{a+h}
\||\nabla|^{m-N}\ \|_p \leq \sum_{m=N}^{N+h}
\||\nabla|^{m-N} \ \|_p = \sum_{k=0}^h
\||\nabla|^k \ \|_p \,.
\]
One has also
\[
\sum_{m=a}^{b+h} \|\ \|_{|\nabla|, N-m, p} = \sum_{m=a}^{b+h}
\||\nabla|^{m-N}\ \|_p \geq \sum_{m=N}^{N+h}
\||\nabla|^{m-N} \ \|_p = \sum_{k=0}^h
\||\nabla|^k \ \|_p \,.
\]


%\smallskip

\eqref{6.1.2}\pointir Since $|\nabla|^{-\mu}$ is a homogeneous pseudodifferential operator of order $-\mu$ and $\mu<Q$, its kernel is homogeneous of degree $-Q+\mu$. According to~\cite[Proposition~1.11]{Folland_1975}, $|\nabla|^{-\mu}$ is bounded from $L^{p}$ to $L^{q}$ if $1<p<q < +\infty$ satisfy $\frac{1}{p} - \frac{1}{q} = \frac{\mu}{Q}$, giving the Hardy--Littlewood--Sobolev inequality
\begin{equation}\label{eq:12}
\||\nabla|^{-\mu}\alpha\|_q \leq C \|\alpha \|_p \,.
\end{equation}
Then for $\alpha \in \Omega_{[a,b]}\cap \mathcal{S}$
\begin{align*}
\sum_{m=a+ \mu}^{b+\mu + h} \| \alpha \|_{|\nabla|, N-m,p} & \geq \sum_{m=N+\mu}^{N+\mu+ h} \||\nabla|^{m-N} \alpha \|_p = \sum_{k=0}^h \| |\nabla|^{\mu + k} \alpha\|_p\\
& \geq 1/C \sum_{k=0}^h \||\nabla|^k \alpha\|_q \asymp
\|\alpha\|_{W^{h,p}_c}
\end{align*}
using~\eqref{eq:12}, \eqref{eq:11} and $|\nabla|^{\mu + k} = |\nabla|^\mu |\nabla|^k $ on $\mathcal{S}$ for $\mu,k\in \N$ as comes from~\cite[Theorem~3.15]{Folland_1975}.
\end{proofc}


\section{Density of \texorpdfstring{$\calS_0$}{S0} and extension of \texorpdfstring{$K_c$}{Kc}}\label{sec:density-s_0}

One knows by Proposition~\ref{prop:K_c_scale} that $K_c$ is continuous with respect to the whole shifted norms $\|\ \|_{|\nabla|, N-m,p}$, but it is only defined and provides an homotopy for $\dd_c$ on the initial domain $\calS_0$ so far. Hence, we have to show that $\calS_0$ is dense in $\calS$ for the standard Sobolev norms $\|\ \|_{W^{h,p}_c}$ in order to extend $K_c$ on forms coming from $\ell^p$ cocycles in Proposition~\ref{prop:leray}. Recall from~\eqref{eq:9} that $\calS_0$ is the space of Schwartz functions with all vanishing polynomial moments. It does not contain any non vanishing function with compact support as seen using Fourier transform or Stone--Weierstrass approximation theorem.

\begin{prop}\label{prop:density-s_0}
For every $h\in\N$ and $1<p<+\infty$, $\calS_0$ is dense in $W^{h,p}_c$. If $p=1$, $\calS_0$ is dense in $\{f\in W_c^{h,1}\,;\,\langle f,1 \rangle=0\}$. If $p=\infty$, $\calS_0$ is dense in the space $C^h_0(G)$ of functions of class $C^h$ that tend to $0$ at infinity.
\end{prop}

\begin{proof}
The group exponential map $\exp:\mathfrak{g}\to G$ maps Schwartz space to Schwartz space, polynomials to polynomials, and Lebesgue measure to Haar measure, hence $\calS_0$ to $\calS_0$. Pick coordinates adapted to the splitting $\mathfrak{g}=\mathfrak{g}_1\oplus\dots\oplus\mathfrak{g}_s$.

First we construct (on $\mathfrak{g}$) a family of functions $(g_\alpha)_{\alpha\in\N^n}$ which is dual to the monomial basis $(x^\beta)_{\beta\in\N^n}$ in the sense that
\[
\langle g_\alpha,x^\beta\rangle=
\begin{cases}
1 & \text{if }\alpha=\beta, \\
0 & \text{otherwise}.
\end{cases}
\]
Let $\F$ denote the Euclidean Fourier transform on $\mathfrak{g}$, let $\Fi$ denote its inverse. Let $n=\dim(G)$, let $\alpha\in\N^n$ denote a multiindex, $|\alpha|=\sum_{i=1}^n \alpha_i$ its usual length, $\alpha!=\prod_{j=1}^n \alpha_j!$. Fix a smooth function $\chi$ with compact support in the square $\{\max |\xi_j| \leq 1\}$ of $\mathfrak{g}^*$, which is equal to $1$ in a neighborhood of $0$. Let
\[
g_\alpha=\Fi\left(\chi(\xi)\frac{(i\xi)^{\alpha}}{\alpha!}\right).
\]
Then $g_\alpha\in\calS$. One computes
\[
(ix)^\beta g_\alpha(x)= \Fi\left(\frac{\partial^{|\beta|}}{\partial\xi^{\beta}} \left(\chi(\xi) \frac{(i\xi)^\alpha}{\alpha!}\right)\right).
\]
Hence
\[
\langle g_\alpha,x^\beta\rangle=(-i)^{|\beta|}\left(\frac{\partial^{|\beta|}}{\partial\xi^{\beta}} \left(\chi(\xi)\frac{(i\xi)^\alpha}{\alpha!}\right) \right)(0) = (-i)^{|\beta|}\left(\frac{\partial^{|\beta|}}{\partial\xi^{\beta}} \left(\frac{(i\xi)^\alpha}{\alpha!}\right)\right)(0)
\]
which vanishes unless $\beta=\alpha$, in which case it is equal to $1$.

From
\[
\frac{\partial g_{\alpha}}{\partial x^{\gamma}}=(-1)^{|\gamma|}\frac{(\alpha+\gamma)!}{\alpha!} g_{\alpha+\gamma},
\]
it follows that
\[
(ix)^\beta \frac{\partial g_{\alpha}}{\partial x^{\gamma}} = (-1)^{|\gamma|}\frac{(\alpha+\gamma)!}{\alpha!}\Fi \left(\frac{\partial^{|\beta|}}{\partial \xi^\beta} \left(\chi(\xi)\frac{(i\xi)^{\alpha+\gamma}}{(\alpha+\gamma)!}\right) \right).
\]
Hence
\[
\left\|x^\beta \frac{\partial g_{\alpha}}{\partial x^{\gamma}}
\right\|_{\infty}\leq \frac{1}{\alpha!} \left\|\frac{\partial^{|\beta|}}{\partial
\xi^\beta} (\chi(\xi)\xi^{\alpha+\gamma})\right\|_1.
\]
Thanks to Leibnitz' formula,
\begin{align*}
\frac{\partial^{|\beta|}}{\partial \xi^\beta} (\chi(\xi)
\xi^{\alpha+\gamma}) &=\sum_{\eta+\eta'=\beta}
\binom{\beta}{\eta}
\frac{\partial^{|\eta|}\chi}{\partial \xi^\eta}
\frac{\partial^{|\eta'|}\xi^{\alpha +\gamma}}{\partial \xi^{\eta'}}\\
&=\sum_{\eta+\eta'=\beta}
\binom{\beta}{\eta}
\frac{\partial^{|\eta|}\chi}{\partial \xi^\eta}
\frac{(\alpha+\gamma)!}{(\alpha+\gamma-\eta')!} \xi^{\alpha+\gamma-\eta'},
\end{align*}
where $\binom{\beta}{\eta}=\prod_{j=1}^n\binom{\beta_j}{\eta_j}$ and, by convention, $\xi^{\alpha+\gamma-\eta'}=0$ if $\alpha +\gamma -\eta'\notin\N^n$. If $\beta$ and $\gamma$ are fixed, this is a polynomial in $\alpha$ of bounded degree $|\beta|$ and bounded coefficients. So is its $L^1$ norm. Therefore
\[
\left\|x^\beta \frac{\partial g_{\alpha}}{\partial x^{\gamma}}\right\|_{\infty}\leq \frac{P_{\beta,\gamma}(\alpha)}{\alpha!},
\]
where $P_{\beta,\gamma}$ is a polynomial on $\R^n$.

Let $w(j)$ denote the weight of the $j$-th basis vector. For $\beta\in\N^n$, let $w(\beta)=\sum_{j=1}^n \beta_j w(j)$. Let $g_{\alpha,t}=t^{w(\alpha)+Q}g_\alpha\circ\delta_t$. Then, for all $\beta$, $\langle x^\beta,g_{\alpha,t}\rangle=t^{w(\alpha)-w(\beta)}\langle x^\beta,g_\alpha\rangle$. Hence, for every choice of sequence $(t_\alpha)$, the family $(g_{\alpha,t_\alpha})_{\alpha\in\N^n}$ is again dual to the monomial basis. Also,
\[
x^\beta \frac{\partial g_{\alpha,t}}{\partial x^{\gamma}} = t^{w(\alpha)-w(\beta)+w(\gamma)+Q} \left(x^\beta \frac{\partial g_{\alpha}}{\partial x^{\gamma}}\right) \circ \delta_t,
\]
hence, for every $1\leq p\leq \infty$ and $p'=\frac{p}{p-1}$,
\[
\left\|x^\beta \frac{\partial g_{\alpha,t}}{\partial x^{\gamma}}\right\|_{p} = t^{w(\alpha)-w(\beta)+w(\gamma)+\frac{Q}{p'}} \left\|x^\beta \frac{\partial g_{\alpha}}{\partial x^{\gamma}} \right\|_{p}.
\]

Given an arbitrary sequence $\mathbf{m}=(m_\alpha)_{\alpha\in\N^n}$, and a positive sequence $\mathbf{t}=(t_\alpha)_{\alpha\in\N^n}$, define the series
\[
f_{\mathbf{m},\mathbf{t}}=\sum_{\alpha\in\N^n}m_\alpha g_{\alpha,t_\alpha}.
\]
If $t_\alpha$ and $|m_\alpha| t_\alpha^{w(\alpha)/2}$ stay bounded by $1$, then for every constant $C$, $|m_\alpha| t_\alpha^{w(\alpha)-C}$ stays bounded, the series converges in $\calS$. Indeed, for every $\beta,\gamma\in\N^n$, the sum of $L^\infty$ norms $\|x^\beta \frac{\partial }{\partial x^{\gamma}}m_\alpha g_{\alpha,t_\alpha}\|_{\infty}$ is bounded above by
\[
\sum_{\alpha\in\N^n}|m_\alpha| t_\alpha^{w(\alpha)-w(\beta)+w(\gamma)+Q}\,\frac{P_{\beta,\gamma} (\alpha)}{\alpha!} <\infty.
\]
By construction, all functions $f_{\mathbf{m},\mathbf{t}}$ have prescribed moments $\langle x^\alpha,f_{\mathbf{m},\mathbf{t}}\rangle=m_\alpha$.

Fix a finite set $F$ of pairs $(\beta,\gamma)$ such that $w(\beta)\leq w(\gamma)$. Let $W^{F,p}$ denote the completion of smooth compactly supported functions for the norm
\[
\|f\|_{W^{F,p}}=\max\left\{\left\|x^\beta \frac{\partial }{\partial x^{\gamma}}m_\alpha g_{\alpha,t_\alpha} \right\|_{p}\,;\,(\beta,\gamma)\in F\right\}.
\]
Denote by
\[
N^{F,p}_\alpha :=\max\left\{\left\|x^\beta \frac{\partial g_{\alpha}}{\partial x^{\gamma}}\right\|_{p} \,;\,(\beta,\gamma)\in F\right\}.
\]
Pick $\mathbf{t}$ such that, in addition to the previous assumptions, $\sum|m_\alpha|t_\alpha^{w(\alpha)+Q/p'}N^{F,p}_\alpha$ converges. Then for every $\eps>0$, the series $f_{\mathbf{m},\eps\mathbf{t}}$ converges in $W^{F,p}$ and for $\eps \leq 1$
\[
\|f_{\mathbf{m},\eps \mathbf{t}}\|_{W^{F,p}}\leq
\eps^{Q/p'}|m_0|N^{F,p}_0 + \eps \sum_{\alpha\not=0} |m_\alpha|t_\alpha^{w(\alpha)+Q/p'} N^{F,p}_\alpha.
\]
Therefore, as $\eps$ tends to $0$, $f_{\mathbf{m},\eps\mathbf{t}}$ tends to $0$ in $W^{F,p}$ (if $p=1$, one must assume that $m_0=0$).

Given $f\in W^{F,p}$ (assume furthermore that $\langle f,1 \rangle=0$ if $p=1$ or that $f\in C_0^h(G)$ if $p=\infty$), approximate $f$ with an element $g\in\calS$ (resp. such that $\langle g,1 \rangle=0$ if $p=1$). Set $m_\alpha=\langle x^\alpha,g \rangle$. Pick $\mathbf{t}$ satisfying the above smallness assumptions with respect to $\mathbf{m}$. Then $g-f_{\mathbf{m},\eps\mathbf{t}}\in \calS_0$ and $f_{\mathbf{m},\eps\mathbf{t}}$ tends to $0$ in $W^{F,p}$, thus $f$ belongs to the closure of $\calS_0$.

Finally, $\|f\|_{W_c^{h,p}}\leq \|f\|_{W^{F,p}}$ for a suitable finite set $F$. Indeed, $G$ admits a basis of left-invariant vectorfields $X_i$ of the form
\begin{equation}\label{eq:13}
X_i=\frac{\partial}{\partial x_i}+\sum_{j>i}P_{i,j}
\frac{\partial}{\partial x_j},
\end{equation}
where $P_{i,j}$ is a $\delta_t$-homogeneous of weight $w(P_{i,j})<w(j)$.
\end{proof}

One can now extend $K_c$ from $\calS_0$ to some Sobolev spaces, depending on the weights on the source and the target. Let $W^{h,p}_c$ denotes the completion of $\calS$, and therefore $\calS_0$, with respect to the norm $\|\ \|_{W^{h,p}_c}$.

\begin{coro}\label{cor:K_c_Sobolev}
Let $a\leq b\leq a+h$ and $a'\leq b'\leq a'+h'$ be integers.
%
Suppose moreover that $1< p<q<\infty$ satisfy $\frac{1}{p} - \frac{1}{q} = \frac{\mu}{Q}$ with $0\leq \mu = b-a' < Q$ and $h= h' + b'-a' + b-a$.

Let $(K_c)_{[a',b']}$ denotes the components of $K_c$ lying in $\Omega_{[a',b']}$. Then $K_c$ extends continuously on $\Omega_{[a,b]} \cap W^{h,p}_c$ so that
\[
(K_c)_{[a',b']}(\Omega_{[a,b]} \cap W^{h,p}_c) \subset W^{h',q}_c \,.
\]
\end{coro}

\begin{proof}
Let $\alpha \in \Omega_{[a,b]}\cap \calS_0$. Apply first Proposition~\ref{prop:Sobolev}$\MK$\eqref{6.1.2} to $(K_c)_{[a',b']}(\alpha)$ with $a'$, $b'$, $\mu$ and $h'$ as above
\begin{align*}
C\|(K_c)_{[a',b']} (\alpha)\|_{W^{h',q}_c} & \leq \sum_{m= a' +
\mu}^{b'+\mu+ h'} \|K_c
\alpha\|_{|\nabla|, N-m,p} = \sum_{m= b}^{a+h} \|K_c
\alpha\|_{|\nabla|, N-m,p} \\
& \leq C' \sum_{m= b}^{a+h} \|\alpha\|_{|\nabla|, N-m,p} \ \text{ by Proposition~\ref{prop:K_c_scale},}\\
& \leq C'' \|\alpha \|_{W^{h,p}_c}\phantom{\sum^{a}_a}
\end{align*}
by Proposition~\ref{prop:Sobolev}$\MK$\eqref{6.1.1} since $\alpha \in \Omega_{[a,b]}$.
\end{proof}

\begin{rema}\label{rem:K_c_formes_lisses}
The statement becomes simpler when forms are smooth, meaning in $W^{+\infty,p}$, as those coming from $\ell^p $-cocycles on $G$. Namely in that case, one can let $h, h'$ go to $+\infty$ with $b'=Q$, $a=0$. Then $K_c$ integrates smooth forms in $W^{+\infty,p}$ and of weight lower than $b$, into smooth forms whose components of weight larger than $a'$ lies in $W^{+\infty,q}$. Note that the $(p,q)$ gap is only determined by the weight gap $\mu = b-a'$ using Sobolev rule.
\end{rema}

\section{Proof of Theorem~\ref{1}\texorpdfstring{$\MK$}{}(\ref{1.3.1})}

Fix a left-invariant Riemannian metric $\mu$ on $G$. Let $\nabla^\mu$ denote its Levi-Civita connection. Let $\nabla^\mathfrak{g}$ denote the connection which makes left-invariant vectorfields parallel. Since $\nabla^\mu-\nabla^\mathfrak{g}$ is left-invariant, hence $\nabla^\mathfrak{g}$-parallel, higher covariant derivatives computed using either $\nabla^\mathfrak{g}$ or $\nabla^\mu$ determine each other via bounded expressions. Therefore the Riemannian Sobolev space $W^{h,p}$ can be alternatively defined using $\nabla^\mathfrak{g}$, i.e. using derivatives along left-invariant vectorfields. Since every left-invariant vectorfield is a combination of compositions of at most $r$ horizontal derivatives $X_i$, where $r$ is the step of $\mathfrak{g}$,
\begin{equation}\label{eq:14}
W^{h,p}\subset W_c^{h,p}\subset W^{h/r,p}.
\end{equation}

According to Proposition~\ref{prop:leray}, it is sufficient to find some large integer $h$ such that every (usual) closed differential form $\alpha \in W^{h,p}\Omega^k(G)$ writes $\alpha = \dd \beta$ with $\beta \in W^{H,q}\Omega^{k-1}(G)$ for $H=n(n+1)$. The relevant value of $h$ will arise from the proof.

We first retract $\alpha$ in the sub-complex $(E,\dd)$ using the differential homotopy $\Pi_E$. Indeed from Theorem~\ref{thm:dc-complex},
\[
\alpha = \Pi_E \alpha + \dd R\alpha + R \dalpha = \Pi_E \alpha + \dd R
\alpha \,,
\]
where $\Pi_E$ and $R$ are left invariant differential operators of (horizontal) order at most $Q$, the homogeneous dimension of $\mathfrak{g}$. Then $\Pi_E \alpha$, $R\alpha$ and $\dd R \alpha$ all belong to $W^{h-Q,p}_c$, so that $\alpha$ and $\Pi_E \alpha$ are homotopic in $L^{p,p}$, and \emph{a fortiori} $L^{q,p}$ Sobolev cohomology for $q\geq p$ (see Remark~\ref{rem:decreasing}).

We now deal with $\alpha_E= \Pi_E \alpha$. Its algebraic projection on $E_0$, $\alpha_c = \Pi_{E_0} \alpha_E$ is a contracted $\dd_c$-closed differential form in $W_c^{h-Q,p}$ too. Since $\alpha_c \in E_0^k \simeq H^k(\mathfrak{g})$, the weights of its components belong to the interval $[a,b]$ where $a=w_{\min} (k)$ and $b=w_{\max} (k)$. Moreover, since $K_c \alpha_c \in E_0^{k-1}$, it belongs to $\Omega_{[a',b']}$ with $a' = w_{\min}(k-1)$ and $b' = w_{\max}(k-1)$.

We can now apply Corollary~\ref{cor:K_c_Sobolev}, with
\[
\mu = b- a' = w_{\max}(k) - w_{\min}(k-1) = \delta N_{\max}(k)
\]
and $h' = h-Q + a-b + a' -b'$. (Observe that $\mu < Q$ except for $G=\R$ in which case $L^1$ one forms have bounded primitives.) We get that $\beta_c = K_c \alpha_c \in W^{h',q}_c$ with $\dd_c \beta_c = \alpha_c$ and $\frac{1}{p}-\frac{1}{q}=\frac{\delta N_{\max}(k)}{Q}$.

Finally, let $ \beta_E=\Pi_E \beta_c$. Then $\beta_E\in W_c^{h'-Q,q}$ since $\Pi_E$ is a differential operator of order at most $Q$. By construction, $\alpha_E=d\beta_E$. By inclusions~\eqref{eq:14}, $\beta_E$ belongs to the Riemannian Sobolev space $W^{(h'-Q)/r,q}$. Thus let us choose $h=rH + 2Q+b-a+b'-a'$. We have shown that every $k$-form in $W^{h,p}$ is homotopic in $L^{p,p}$ Sobolev cohomology to a form that has a primitive in $W^{H,q}$, where
\[
\frac{1}{p}-\frac{1}{q}=\frac{\mu}{Q}= \frac{\delta N_{\max}(k)}{Q}.
\]
Proposition~\ref{prop:leray} implies that $\ell^{q,p}H^k(G)=0$. \emph{A fortiori}, $\ell^{q',p}H^k(G)=0$ for all $q'\geq q$.

\subsection{Remarks on the weight gaps}\label{sec:remarks-weight-gaps}

Implicit in the statement of Theorem~\ref{1} is that for any $1\leq k \leq n $, the weight gap $\delta N(k) = w_{\max}(k)- w_{\min}(k-1)$ is positive. Actually, one has
\begin{equation}\label{eq:15}
w_{\max}(k) - w_{\max}(k-1) \geq 1 \quad \mathrm{and}\quad w_{\min}(k) - w_{\min}(k-1) \geq 1.
\end{equation}

\begin{proof}
Since $E_0 = \ker \dd_0 \cap \ker \delta_0$ is Hodge-$*$ symmetric with $w(*\alpha) = Q - w (\alpha)$ (see e.g.~\cite{Rumin_CRAS_1999}), one can restrain to prove the statement about $w_{\min}$. Let $\alpha \in H^{k}(\mathfrak{g})$ be non zero with minimal weight $w_{\min}(k)$. See it as a left invariant retracted form in $E_0^k$. Following Section~\ref{thm:dc-complex}, one has then $\dalpha = \dd_0 \alpha = \dd_c \alpha = 0$. Now $(E_0,\dd_c)$ being locally exact (homotopic to de Rham complex), one has $\alpha = \dd_c \beta$ for some $\beta \in E_0^{k-1}$. Since $\dd_c$ increases the weight by $1$ at least (as $\dd_0=0$ on it), $\beta$ has a non vanishing component of weight $< w_{\min}(k)$, whence $w_{\min}(k-1) \leq w_{\min}(k) -1$.
\end{proof}

Another remark is about the use of the contracted complex here. As observed in Remark~\ref{rem:1}, de Rham complex being C-C elliptic too, one could directly use a similar homotopy $K$ for it in the previous proof. But then, it would lead to a weaker integration result of closed $W_c^p$ forms in $W_c^q$ for a larger gap $\frac{1}{p} - \frac{1}{q} = \frac{\delta N}{Q}$. Indeed this $\delta N $ is the maximal weight gap between the whole $\Omega^k(G)$ and $\Omega^{k-1}(G)$, instead of the restricted ones. This gives a weaker vanishing condition of $\ell^{q,p} H^{k}(G)$. The point here is that the first homotopy $\Pi_E$ to $E$, being differential, ``costs a lot of derivatives'' that would be a shame locally, but is harmless here when working with the quite smooth $W_c^{h,p}$ forms coming from our $\ell^p$ simplicial cocycles. Indeed, we have seen that $\Pi_E = 1$ in $L^{p,p}$ Sobolev cohomology.

%\smallskip

Still about these ideas of ``loosing'' or ``gaining'' derivatives, things go exactly in the opposite direction as usual here. Namely, one sees in the proof that the higher derivatives the homotopy $K_c$ controls at some place, the larger becomes the $(p,q)$ gap from Hardy--Littlewood--Sobolev inequality in Corollary~\ref{cor:K_c_Sobolev}. This is of course better in local problems since $L^q_{loc}$ gets smaller, but is weaker on global smooth forms as $W_c^q$ gets larger. In this large scale low frequency integration problem, the less derivatives you gain is the better. No gain, no pain.


\section{Nonvanishing result}\label{nonvanishing}

As we will see, in order to prove the non-vanishing of the $\ell^{q,p}$ and Sobolev $L^{q,p}$ cohomology of $G$ for some $1\leq p \leq q\leq +\infty$, we shall construct closed $k$-forms $\omega\in W^{h,p}(G)$, $h$ large enough, and $n-k$-forms $\omega'_j$ such that $\|\dd\omega'_j\|_{q'}$ tend to zero whereas $\int_{G}\omega\wedge\omega'_j$ stays bounded away from $0$. Here $q'$ is the dual exponent of $q$ : $\frac{1}{q}+\frac{1}{q'}=1$.


The building blocks will be differential forms which are homogeneous under dilations $\delta_t$. In this section, one can use any expanding one-parameter group $s\mapsto h_s$ of automorphisms of $G$. Expanding means that the derivation $D$ generating $h_s$ has positive eigenvalues. The one-parameter group $s\mapsto \delta_{e^s}$ is an example, but others may be useful, see below.

\subsection{Homogeneous differential forms}\label{sec:homog-diff-forms}

Fix an expanding one-parameter group $s\mapsto h_s$ of automorphisms of $G$, generated by a derivation $D$. The data $(G,(h_s))$ is called a \emph{homogeneous Lie group}. Denote by $T=\trace(D)$ its
\emph{homogeneous dimension}.

Left-invariant differential forms on $G$ split into weights $w$ under $h_s$.

Say a smooth differential form $\omega$ on $G\setminus\{1\}$ is homogeneous of degree $\lambda$ if $h_s^*\omega=e^{s\lambda} \omega$ for all $s\in\R$. Note that homogeneity is preserved by $\dd$. A left-invariant differential form of weight $w$ is homogeneous of degree $w$

Let $\rho$ denote a continuous function on $G$ which is homogeneous of degree~$1$, and smooth and positive away from the origin. Let $\beta$ be a nonzero continuous differential form which is homogeneous of degree $\lambda$, smooth away from the origin and has weight $w$, then
\[
\beta=\rho^{\lambda-w}\sum a_i\theta_i,
\]
where $\theta_i$'s are left-invariant of weight $w$ and $a_i$ are smooth homogeneous functions of degree~$0$, hence are bounded. Therefore
\[
\beta\in L^p(\{\rho\geq 1\}) \iff
\int_{1}^{+\infty}\rho^{p(\lambda-w)}\rho^{T-1}\, \drho<\infty
\iff \lambda-w +\frac{T}{p}<0.
\]
It follows that if $\gamma$ is a differential form which is homogeneous of degree $\lambda$ and has weight $\geq w$,
\begin{equation}\label{eq:16}
\lambda-w +\frac{T}{p}<0 \implies \gamma\in L^p(\{\rho\geq 1\}).
\end{equation}

Start with a closed differential $k$-form $\omega$ which is homogeneous of degree $\lambda$ and of weight $\geq w$ (and no better). Pick a differential $n-k$-form $\alpha'$ which is homogeneous of degree $\lambda'$. Assume that $\dalpha'$ has weight $\geq w'$ (and no better). Set
\[
\omega'_j=\chi_j \alpha',
\]
where $\chi_j=\chi\circ\rho\circ h_{-j}$ and $\chi$ is a cut-off supported on $[1,2]$. The top degree form $\omega\wedge \alpha'$ is homogeneous of degree $\lambda+\lambda'$ and has weight $T$. It belongs to $L^1$ if $\lambda+\lambda'<0$. Thus, in order that $\int\omega\wedge \omega'_j$ does not tend to $0$, it is necessary that $\lambda+\lambda'\geq 0$.

By construction,
\[
\omega'_j=e^{\lambda'j} h_{-j}^*\omega'_1.
\]
Let $\beta$ denote a component of $\dd\omega'_1$ of weight $\tilde w$, and $s=-j$. By the change of variable formula,
\begin{align*}
\| h_s^*\beta\|_{q'}^{q'} &= \int | h_s^*\beta|^{q'}\,\dvol\\
&= \int e^{q'{\tilde w}s}|\beta|^{q'}\circ h_s\,\dvol\\
&= e^{(q'{\tilde w}-T)s}\int |\beta|^{q'}\,\dvol\\
&= e^{(q'{\tilde w}-T)s}\|\beta\|_{q'}^{q'}.
\end{align*}
This works as well for $q'=\infty$. Since $d\omega'_1$ has weight $\geq w'$,
\[
\|\omega'_j\|_{q'}\leq e^{(\lambda'-w'+\frac{T}{q'})j}\|\omega'_1\|_{q'}.
\]
One concludes that
\[
\|\omega'_j\|_{q'}\to 0 \iff \lambda'-w'+\frac{T}{q'}<0.
\]
It holds for $q'=\infty$ as well. Remember that $\omega\in L^p \iff \lambda-w+\frac{T}{p}<0$. Note that
\[
\lambda-w+\frac{T}{p}+\lambda'-w'+\frac{T}{q'} =\lambda+\lambda'-w-w'+T+T\left(\frac{1}{p}-\frac{1}{q}\right).
\]
Finally, one sees that one can pick $\lambda$ and $\lambda'$ such that $\omega\in L^p$, $\|\omega'_j\|_{q'}\to 0$ and $\lambda+\lambda'\geq 0$ iff
\begin{equation}\label{eq:17}
\frac{1}{p}-\frac{1}{q}<\frac{w+w'-T}{T}.
\end{equation}

\subsection{A numerical invariant of homogeneous groups}

A lower bound for the sum $w+w'$ appearing in above inequation~\eqref{eq:17} is provided by the following definitions.

Let $\Sigma$ denote the level set $\{\rho=1\}$. It is a smooth compact hypersurface, transverse to the vectorfield $\xi$ which generates the 1-parameter group $s\mapsto h_{e^s}$. Differential $k$-forms which are homogeneous of degree $\lambda$ on $G\setminus\{1\}$ correspond to smooth sections of the pull-back of the bundle $\Lambda^k T^* G$ by the injection $\Sigma\hookrightarrow G$. Given such a section $\sigma$, a form $\alpha$ is defined as follows. At a point $x$ where $\rho(x)=e^s$, $\alpha(x)= h_{-s}^*\sigma$. Conversely, given a homogeneous form $\alpha$, consider its values $\sigma$ along $\Sigma$ (not to be confused with the restriction of $\alpha$, which belongs to $\Lambda^k T^* \Sigma$). A similar construction applies to contracted forms as $E_0$ is stable by dilations.

On spaces of homogeneous forms of complementary degrees $k$ and $n-k$ and complementary degrees of homogeneity $\lambda$ and $-\lambda$, define a pairing as follows: if $\beta$ and $\beta'$ are homogeneous of degrees $\lambda$ and $-\lambda$, set
\[
I(\beta,\beta')=\int_{\Sigma}\iota_\xi(\beta\wedge\beta').
\]
This is a nondegenerate pairing. Indeed, pointwise, an $n$-form $\omega$ is determined by the restriction of $\iota_\xi(\omega)$ to $T\Sigma$, hence the pointwise pairing $(\beta,\beta')\mapsto \iota_\xi(\beta\wedge\beta')_{|T\Sigma}$ is nondegenerate. For instance, one has $\beta \wedge * \beta = \|\beta\|^2 \,\dvol$ pointwise. Note that the $n$-form $\beta\wedge\beta'$ is homogeneous of degree~$0$, i.e. dilation invariant. The $n-1$-form $\iota_\xi(\beta\wedge\beta')$ is closed, the integral $I(\beta,\beta')$ only depends on the cohomology class of this form. The boundary of any smooth bounded domain containing the origin can be used to perform integration instead of $\Sigma$.


\begin{defi}\label{dfi:wsG}
Let $G$ be a homogeneous Lie group of homogeneous dimension $T$. For $k=1,\dots,n= \dim(G)$, define $ws_G(k)$ as the maximum of sums $w+w'-T$ such that for a dense set of real numbers $\lambda$, there exist
\begin{enumerate}
\item a differential form $\alpha$ of degree $k-1$ on $G\setminus\{1\}$, homogeneous of degree $\lambda$, such that $\dalpha$ has weight $\geq w$,
\item a differential form $\alpha'$ of degree $n-k$ on $G\setminus\{1\}$, homogeneous of degree $\lambda'=-\lambda$, such that $\dalpha'$ has weight $\geq w'$,
\end{enumerate}
such that $ I(\dalpha,\alpha')\not=0$.
\end{defi}

Note that for all $k\geq 1$, $ws_G(k)=ws_G(n-k+1)$. For instance, when Carnot dilations are used, nonzero $1$-forms of weight $\geq 2$ are never closed, and $n$-forms are always closed and of weight $T$, hence $ws_G(1)=ws_G(n)=1$.

\subsection{Cohomology nonvanishing}

\begin{theo}\label{thm:nonzero}
Let $G$ be a homogeneous Lie group of homogeneous dimension $T$. Then $\ell^{q,p}H^k(G)\not=0$ provided
\[
1\leq p,q<+\infty \quad \text{and} \quad
\frac{1}{p}-\frac{1}{q}<\frac{ws_G(k)}{T}.
\]
\end{theo}

\begin{proof}
By assumption, $\epsilon=\frac{w+w'-T}{T}-\frac{1}{p}+\frac{1}{q}>0$. Pick a real number $\lambda$ in the dense set given in Definition~\ref{dfi:wsG} and such that
\begin{equation}\label{eq:18}
w-\frac{T}{p}-\frac{T}{2}\epsilon \leq \lambda < w-\frac{T}{p}.
\end{equation}
Then $\lambda'=-\lambda$ satisfies
\begin{align}
\lambda'-w'+\frac{T}{q'} & = -\lambda - w' + \frac{T}{q'} \nonumber \\
& \leq -w + \frac{T}{p} + \frac{T\epsilon}{2} -w' + \frac{T}{q'} = - \frac{T\epsilon}{2} < 0. \label{eq:19}
\end{align}
By definition, there exist differential $k-1$ and $n-k$-forms $\alpha$ and $\alpha'$, homogeneous of degrees $\lambda$ and $\lambda'$, such that $\dalpha$ and $\dalpha'$ have weights $\geq w$ and $\geq w'$. Then $\dalpha \wedge \alpha'$ is homogeneous of degree~$0$. Using the notations of Section~\ref{sec:homog-diff-forms}, for all $j$, $\dalpha\wedge \chi_j\alpha'= h_{-j}^*(\dalpha\wedge \chi_1
\alpha')$, hence
\[
\int_{G} \dalpha \wedge \chi_j\alpha'=\int_{G} \dalpha \wedge
\chi_1\alpha'=I(\dalpha,\alpha')\int_{\R}\chi(t)\,\dt\not=0.
\]
Since $\dalpha$ is homogeneous of degree $\lambda$ and has weight $\geq w$, it belongs to $L^p$ (away from a neighborhood of the origin) by~\eqref{eq:16} and~\eqref{eq:18}. Furthermore, derivatives along left invariant vector fields decrease homogeneity. Hence all such derivatives of $\dalpha$ belong to $L^p$. After smoothing $\alpha$ near the origin, we get a closed form $\omega$ on $G$ that coincides with $\dalpha$ on $\{\rho \geq 1\}$, and which belongs to $W^{h,p}$ for all $h$. Set
\[
\omega'_j=\chi_j \alpha'.
\]
Then $\int_{G}\omega\wedge \omega'_j$ does not depend on $j$.

Assume by contradiction that $\omega=\dphi$ where $\phi\in W^{h,q}$. In particular, $\phi\in L^q$. Since $\omega'_j$ are compactly supported, Stokes theorem applies and
\begin{align*}
\left|\int_{G}\omega\wedge \omega'_j\right| &=\left|\int_{G}\dphi\wedge \omega'_j\right|\\
&=\left|\int_{G}\phi\wedge \domega'_j\right|\\
&\leq\|\phi\|_{q}\|\domega'_j\|_{q'}
\end{align*}
which tends to $0$, as $\alpha'$ and $\dalpha' \in L^q(\{\rho \geq 1\}$) by~\eqref{eq:16} and~\eqref{eq:19}, contradiction. We conclude that $[\omega]\not=0$ in the Sobolev $L^{q,p}$ cohomology of $G$. According to Proposition~\ref{prop:leray}, this implies that the $\ell^{q,p}$ cohomology of $G$ does not vanish.
\end{proof}

\subsection{Lower bounds on \texorpdfstring{$ws_G$}{wsG}}

We give here two lower bounds on $ws_G$. Combined with Theorem~\ref{thm:nonzero}, they complete the proof of Theorem~\ref{1}$\MK$\eqref{1.3.2} in the wider setting of homogeneous groups. We start with a lemma on the contracted complex.

\begin{lemm}\label{lem:existence_formes}
Let $G$ be a homogeneous Lie group of dimension $n$, let $k=1,\dots,n$. Then for an open dense set of real numbers $\lambda$, there exist smooth non $\dd_c$-closed contracted $k-1$-forms on $G \setminus \{1\}$ which are homogeneous of degree $\lambda$.
\end{lemm}

\begin{proof}
The differential $\dd_c \not\equiv 0$ on $E_0^{k-1}$, since the complex $(\dd_c, E_0^*)$ is a resolution on $G$. Then, by the Stone--Weierstrass approximation theorem, their exist non $\dd_c$-closed contracted forms with homogeneous polynomial components in an invariant basis. Pick one term $P \alpha_0$ with $\alpha_0\in E_0^{k-1}$ invariant, and a non constant homogeneous polynomial $P$ such that $\dd_c (P\alpha_0)\not=0$. Up to changing $P$ into $-P$, pick $x_0 \in G$ such that $\dd_c(P\alpha_0) (x_0) \not=0$ and $P(x_0) >0$. Consider the map
\[
F \ :\ \lambda \in \mathbb{C} \mapsto \dd_c (P^\lambda \alpha_0) (x_0).
\]
Since $\dd_c$ is a differential operator, $F$ is analytic. Since $F(1) \not=0$, one has $F(\lambda) \not= 0$ except for a set of isolated values of $\lambda$. Let $\chi$ be a smooth homogeneous function on $G \setminus \{1\}$ of degree~$0$ with support in $\{P> 0\}$ and $\chi=1$ around $x_0$. Then $\alpha = \chi P^\lambda \alpha_0$ is a smooth non $\dd_c$-closed homogeneous contracted form on $G\setminus \{1\}$ of degree $w(\alpha)= \lambda w(P) + w(\alpha_0)$.
\end{proof}

\begin{prop}\label{lower}
Let $G$ be a homogeneous Lie group of dimension $n$. For all $k=1,\dots,n$,
\[
ws_G(k)\geq \max\{1, w_{\min}(k)-w_{\max}(k-1)\}.
\]
\end{prop}

\begin{proof}
By Lemma~\ref{lem:existence_formes}, pick a non $\dd_c$-closed contracted $k-1$-form $\alpha$, homogeneous of degree $\lambda$. Assume that $\dd_c\alpha$ has weight $\geq w$ and no better (i.e. its weight $w$ component $(\dd_c \alpha)_w$ does not vanish identically). Pick a smooth contracted $n-k$-form $\alpha'$ of weight $T-w$, homogeneous of degree $-\lambda$ and such that $I(\dd_c\alpha,\alpha')\not=0$. For instance $\alpha'= \rho^{-2\lambda+2w-T} * (\dd_c\alpha)_w$ will do. Set $\alpha_E = \Pi_E \alpha$ and $\alpha_E'= \Pi_E \alpha'$.

By construction (see Theorem~\ref{thm:dc-complex}), $\Pi_E = \Pi_{E_0}+ D$ where $D$ strictly increases the weight. Hence $\dalpha_E - \dd_c \alpha= \Pi_E \dd_c \alpha - \dd_c\alpha$ has weight $\geq w +1$, and $\alpha'_E - \alpha'$ has weight $\geq T-w + 1$. Therefore $\dalpha_E \wedge \alpha'_E - \dd_c \alpha \wedge \alpha'$ has weight $\geq T+1$, thus vanishes. Then it holds that
\[
I(\dalpha_E, \alpha'_E) = I(\dd_c \alpha, \alpha') \not=0 \,.
\]
Consider now the weight of $\dalpha'_E$. By construction, $E \subset \ker \dd_0$, so that $\dd$ strictly increases the weight on $E$, see Section~\ref{sec:dc_complex}. Therefore
\[
w(\dalpha'_E) \geq w(\alpha'_E) + 1 = w(\alpha')+ 1 = T - w(\dd_c
\alpha) + 1 = T -w(\dalpha_E) + 1,
\]
hence $ws_G(k) \geq w(\dalpha_E) + w (\dalpha'_E) - T \geq 1 $ as needed. One has also that
\[
w(\dalpha'_E) = w (\dd_c \alpha') \geq w_{\min}(n-k+1) = T - w_{\max}(k-1),
\]
by Hodge $*$-duality, see proof of~\eqref{eq:15}, while
\[
w(\dalpha_E) = w (\dd_c \alpha) \geq w_{\min}(k) \,.
\]
This gives $ws_G(k) \geq w(\dalpha_E) + w (\dalpha'_E) - T \geq w_{\min}(k) - w_{\max}(k-1)$.
\end{proof}


\subsection{An example: Engel's group}\label{sec:engels-group}

We illustrate the non-vanishing results on the Engel group $E^4$.

It has a 4-dimensional Lie algebra with basis $X,Y,Z,T$ and nonzero brackets $[X,Y]=Z$ and $[X,Z]=T$. One finds, see e.g.~\cite[Section~2.3]{Rumin_TSG}, that
\[
H^1(\mathfrak{g}) \simeq \mathrm{span}(\theta_X, \theta_Y) \
\mathrm{and} \ H^2(\mathfrak{g}) \simeq \mathrm{span}(\theta_Y \wedge
\theta_Z, \theta_X \wedge \theta_T) \,.
\]
The following table gives the values of $\delta N_{\max}$ and $\delta N_{\min}$ for $E^4$ with respect to its standard Carnot weight : $w(X)= w(Y) =1$, $w(Z)= 2$ and $w(T) = 3$. One has $Q=7$ and
\[
\begin{tabular}{|l|c|c|c|c|}
\hline k & 1 & 2 & 3 & 4 \\
\hline\hline $w_{\max}(k)$ & 1 & 4 & 6 & 7 \\
\hline $w_{\min}(k)$ & 1 & 3 & 6 & 7 \\
\hline $\delta N_{\max}(k)$ & 1 & 3 & 3 & 1 \\
\hline $\delta N_{\min}(k)$ & 1 & 2 & 2 & 1 \\
\hline
\end{tabular}
\]

We see that Theorem~\ref{1} is sharp in degrees $1$ and $4$. However, there are gaps in degrees $2$ and $3$. In particular, $H^{2,q,p}(E^4)$ vanishes when $\frac{1}{p}- \frac{1}{q} \geq \frac{3}{7}$ and does not when $\frac{1}{p}- \frac{1}{q} < \frac{2}{7}$, provided $1<p, q< +\infty$.

%\smallskip
Following~\cite[Section~4.2]{Rumin_TSG}, let us also use the expanding one-parameter group of automorphisms of $E^4$ generated by the derivation $D$ defined by
\[
D(X)=X,\quad D(Y)=2Y,\quad D(Z)=3Z,\quad D(T)= 4T.
\]
Then $\trace(D)=10$, and with this choice of derivation, the table of weights becomes
\[
\begin{tabular}{|l|c|c|c|c|}
\hline k & 1 & 2 & 3 & 4 \\
\hline\hline $w_{\max}(k)$ & 2 & 5 & 9 & 10 \\
\hline $w_{\min}(k)$ & 1 & 5 & 8 & 10 \\
\hline $\delta N_{\max}(k)$ & 2 & 4 & 4 & 2 \\
\hline $\delta N_{\min}(k)$ & 1 & 3 & 3 & 1 \\
\hline
\end{tabular}
\]

According to Proposition~\ref{lower}, with respect to this homogeneous structure, $ws_{E^4}(2)\geq \delta N_{\min}(2)=3$. Then with Theorem~\ref{thm:nonzero},
\[
1\leq p,q<+\infty \quad \text{and} \quad \frac{1}{p}-
\frac{1}{q}<\frac{3}{10}
\]
implies that $\ell^{q,p}H^2(E^4)\not= \{0\}$. We see that a non-Carnot homogeneous structure may yield a better interval for cohomology nonvanishing, which is intriguing for a large scale geometric invariant on a Carnot group.



\begin{thebibliography}{MMMMMM}
\bibitem[BFP17]{BFP} Baldi, Annalisa; Franchi, Bruno; Pansu, Pierre Poincaré and Sobolev inequalities for differential forms in Heisenberg groups (2017) (\url{https://arxiv.org/abs/1711.09786}) \href{https://ahl.centre-mersenne.org/item/AHL_2018__1__267_0/#BFP}{Publisher reference}.
\bibitem[BK15]{Bourdon-Kleiner} Bourdon, Marc; Kleiner, Bruce Some applications of \(\ell\) p -cohomology to boundaries of Gromov hyperbolic spaces, Groups Geom. Dyn., Volume 9 (2015) no. 2, pp. 435-478 \href{https://ahl.centre-mersenne.org/item/AHL_2018__1__267_0/#Bourdon-Kleiner}{Publisher reference}.
\bibitem[BMOR84]{B-MinOo-Ruh} Bemelmans, Josef; Min-Oo, Maung; Ruh, Ernst A. Smoothing Riemannian metrics, Math. Z., Volume 188 (1984) no. 1, pp. 69-74 \href{https://ahl.centre-mersenne.org/item/AHL_2018__1__267_0/#B-MinOo-Ruh}{Publisher reference}.
\bibitem[BP03]{Bourdon-Pajot} Bourdon, Marc; Pajot, Hervé Cohomologie l p  et espaces de Besov, J. Reine Angew. Math., Volume 558 (2003), pp. 85-108 \href{https://ahl.centre-mersenne.org/item/AHL_2018__1__267_0/#Bourdon-Pajot}{Publisher reference}.
\bibitem[CGGP92]{CGGP} Christ, Michael; Geller, Daryl; Głowacki, Paweł; Polin, Larry Pseudodifferential operators on groups with dilations, Duke Math. J., Volume 68 (1992) no. 1, pp. 31-65 \href{https://ahl.centre-mersenne.org/item/AHL_2018__1__267_0/#CGGP}{Publisher reference}.
\bibitem[CT11]{Cornulier-Tessera} Cornulier, Yves; Tessera, Romain Contracting automorphisms and L p -cohomology in degree one, Ark. Mat., Volume 49 (2011) no. 2, pp. 295-324 \href{https://ahl.centre-mersenne.org/item/AHL_2018__1__267_0/#Cornulier-Tessera}{Publisher reference}.
\bibitem[DM16]{Drutu-Mackay} Drutu, Cornelia; Mackay, John M. Random groups, random graphs and eigenvalues of p-Laplacians (2016) (\url{https://arxiv.org/abs/1607.04130}) \href{https://ahl.centre-mersenne.org/item/AHL_2018__1__267_0/#Drutu-Mackay}{Publisher reference}.
\bibitem[Ele98]{Elek} Elek, Gábor Coarse cohomology and l p -cohomology, K-Theory, Volume 13 (1998) no. 1, pp. 1-22 \href{https://ahl.centre-mersenne.org/item/AHL_2018__1__267_0/#Elek}{Publisher reference}.
\bibitem[Fol75]{Folland_1975} Folland, Gerald B. Subelliptic estimates and function spaces on nilpotent Lie groups, Ark. Mat., Volume 13 (1975) no. 2, pp. 161-207 \href{https://ahl.centre-mersenne.org/item/AHL_2018__1__267_0/#Folland_1975}{Publisher reference}.
\bibitem[FS74]{FollandStein} Folland, Gerald B.; Stein, Elias M. Estimates for the \(\partial\) ¯ b  complex and analysis on the Heisenberg group, Commun. Pure Appl. Math., Volume 27 (1974), pp. 429-522 \href{https://ahl.centre-mersenne.org/item/AHL_2018__1__267_0/#FollandStein}{Publisher reference}.
\bibitem[Gen14]{Genton} Genton, Luc Scaled Alexander-Spanier Cohomology and L q,p  Cohomology for Metric Spaces, Ph. D. Thesis, EPFL, Lausanne (Switzerland) (2014) (Thèse no 6330) \href{https://ahl.centre-mersenne.org/item/AHL_2018__1__267_0/#Genton}{Publisher reference}.
\bibitem[Gro93]{Gromov} Gromov, Mikhael Asymptotic invariants of infinite groups, Geometric group theory, Vol. 2 (Sussex, 1991) (London Mathematical Society Lecture Note Series), Volume 182, Cambridge University Press, 1993, pp. 1-295 \href{https://ahl.centre-mersenne.org/item/AHL_2018__1__267_0/#Gromov}{Publisher reference}.
\bibitem[Pan08]{Pansu4} Pansu, Pierre Cohomologie L p  et pincement, Comment. Math. Helv., Volume 83 (2008) no. 2, pp. 327-357 \href{https://ahl.centre-mersenne.org/item/AHL_2018__1__267_0/#Pansu4}{Publisher reference}.
\bibitem[Pan17]{Pansu2} Pansu, Pierre Cup-products in L q,p -cohomology: discretization and quasi-isometry invariance (2017) (\url{https://arxiv.org/abs/1702.04984}) \href{https://ahl.centre-mersenne.org/item/AHL_2018__1__267_0/#Pansu2}{Publisher reference}.
\bibitem[Rum99]{Rumin_CRAS_1999} Rumin, Michel Differential geometry on C-C spaces and application to the Novikov-Shubin numbers of nilpotent Lie groups, C. R. Math. Acad. Sci. Paris, Volume 329 (1999) no. 11, pp. 985-990 \href{https://ahl.centre-mersenne.org/item/AHL_2018__1__267_0/#Rumin_CRAS_1999}{Publisher reference}.
\bibitem[Rum01]{Rumin_TSG} Rumin, Michel Around heat decay on forms and relations of nilpotent Lie groups, Séminaire de Théorie Spectrale et Géométrie, Année 2000–2001 (Séminaire de Théorie Spectrale et Géométrie), Volume 19, Université de Grenoble I, Institut Fourier, 2001, pp. 123-164 \href{https://ahl.centre-mersenne.org/item/AHL_2018__1__267_0/#Rumin_TSG}{Publisher reference}.
\bibitem[Rum10]{Rumin10} Rumin, Michel Spectral density and Sobolev inequalities for pure and mixed states, Geom. Funct. Anal., Volume 20 (2010) no. 3, pp. 817-844 \href{https://ahl.centre-mersenne.org/item/AHL_2018__1__267_0/#Rumin10}{Publisher reference}.
\bibitem[VSCC92]{VCSC} Varopoulos, Nicholas Th.; Saloff-Coste, Laurent; Coulhon, Thierry Analysis and geometry on groups, Cambridge Tracts in Mathematics, 100, Cambridge University Press, 1992, xii+156 pages \href{https://ahl.centre-mersenne.org/item/AHL_2018__1__267_0/#VCSC}{Publisher reference}.
\end{thebibliography}
\end{document}
